\documentclass[sigconf,anonymous,review]{acmart}

\settopmatter{printacmref=false}
\renewcommand\footnotetextcopyrightpermission[1]{}

\usepackage{xspace}
\usepackage{graphicx}
\usepackage{float}
\usepackage{algorithm}
\usepackage{algpseudocode}
\usepackage{hyperref}
\usepackage{pgfplots}
\pgfplotsset{compat=1.18}
\usepackage{pgfplotstable}
\usepackage{tikz}
\usetikzlibrary{arrows.meta, positioning}

\usepackage{amsmath}

\newtheorem*{remark*}{Remark}

% numbers and claims still to be confirmed by a re-run or a literature check
\newcommand{\pending}[1]{\textcolor{red}{#1}}

% a conditional on one line, for algorithms that would not fit a column otherwise
\algnewcommand{\IfThen}[2]{\State \algorithmicif\ #1\ \algorithmicthen\ #2}

\usepgfplotslibrary{groupplots}

% the names of our programs in the experiments
\newcommand{\wfdo}{WF-DO\xspace}
\newcommand{\wfdomem}{WF-DO-mem\xspace}

\AtBeginDocument{%
\providecommand\BibTeX{{%
Bib\TeX}}}

\setcopyright{acmlicensed}
\copyrightyear{2025}
\acmYear{2025}
\acmDOI{XXXXXXX.XXXXXXX}

\acmConference[ICPP]{International Conference on Parallel Processing}{September 8--11, 2025}{at The Catamaran Resort Hotel San Diego, CA}
\acmISBN{978-1-4503-XXXX-X/2025/06}

\title{Parallel Wait-Free BFS}

\begin{abstract}
Breadth-first search (BFS) is a basic graph algorithm. The fastest parallel BFS implementations for shared memory are direction-optimizing: they expand a search level by level, each level top-down or bottom-up. They are also blocking. Every level or round ends at a barrier or a join, so a thread that is preempted, interrupted or stopped by a page fault while it holds part of the work delays all the others, at every level that the delay overlaps. We present a wait-free method for direction-optimizing BFS, in which threads cooperate on a single search and every call returns after a bounded number of its own steps, whatever the other threads do.

Each level holds one bucket per thread. A thread expands its own bucket, then helps with what is left of the others, and orders its writes so that a thread that stops at any point leaves nothing that the others cannot redo. The threads agree on the direction of each level with a single compare-and-swap, bottom-up levels divide the vertices into blocks that are owned and helped like buckets, and one-byte codes give the bottom-up step an exact test of the current level, even for threads that have fallen behind. We prove the method correct and wait-free under the memory model of x86, and describe a variant whose memory depends on the largest levels of a search rather than on all of them.

On a 128-core, two-socket machine, against GAPBS, GBBS and PASGAL, our method is the fastest on four of five graphs, by 1.33 to 1.89$\times$ over the fastest of the three; on the road network, whose search has 6,262 levels, the memory optimized variant is 1.09$\times$ faster than PASGAL. On LiveJournal, when one of 16 threads stalls for 10~ms at the start of every level or round, GAPBS, GBBS and PASGAL take 2.8 to 9.0$\times$ as long as without stalls, while our method takes 16\% longer.
\end{abstract}


\begin{CCSXML}
<ccs2012>
   <concept>
       <concept_id>10010147.10010169.10010170.10010171</concept_id>
       <concept_desc>Computing methodologies~Shared memory algorithms</concept_desc>
       <concept_significance>500</concept_significance>
   </concept>
</ccs2012>
\end{CCSXML}

\ccsdesc[500]{Computing methodologies~Shared memory algorithms}

\keywords{Breadth-first search (BFS), Direction-optimizing BFS, Wait-free algorithms, Helping, Parallel graph algorithms, Shared memory}



\theoremstyle{definition}
\newtheorem{definition}{Definition}[section]

\begin{document}

\maketitle

\section{Introduction}
\label{sec:intro}

Breadth-first search (BFS) is one of the most basic graph algorithms. Starting from a source vertex, it visits the graph level by level, where level $d$ holds the vertices at distance $d$ from the source, and so finds shortest paths in unweighted graphs. It is a building block of many other graph algorithms, such as connectivity, betweenness centrality and maximum flow, and it is the search kernel of the Graph500 and GAP benchmarks~\cite{graph500, gapbs}.

In most parallel BFS algorithms for shared memory, the threads expand the vertices of one level together, which discovers the next level, and move on to it only once the whole level is done; such algorithms are \emph{level-synchronous}~\cite{bader_bfs, charles_bfs}. Beamer et al.~\cite{beamer_dobfs} observed that on low-diameter graphs, a few middle levels hold most of the vertices, and that on those levels it is cheaper to search \emph{bottom-up}: every unvisited vertex looks for a neighbor in the current level and stops at the first one it finds, instead of every vertex of the level checking all of its neighbors. Their \emph{direction-optimizing} BFS chooses between the two directions at every level. It underlies the fastest shared-memory BFS implementations today, including those of GAPBS~\cite{gapbs}, Ligra~\cite{ligra}, GBBS~\cite{gbbs} and PASGAL~\cite{pasgal}.

These implementations update vertices with atomic instructions rather than locks, but their BFS methods are still \emph{blocking}. GAPBS ends every step with an OpenMP barrier, and GBBS and PASGAL end every round when its parallel loop joins, so a level is over only once its slowest thread has finished its share. A thread that is delayed while it holds part of a level, because the operating system preempted it for another process, or an interrupt or a page fault stopped it, keeps every other thread waiting at the end of that level. Dynamic scheduling and work stealing~\cite{blumofe_workstealing} hand the work that a delayed thread has not started to other threads, but not the work it is in the middle of. The cost recurs at every level that a delay overlaps, an effect well known from large parallel machines, where short delays on single processors, amplified by synchronization, slow down whole applications~\cite{petrini_noise, hoefler_noise}. The number of levels, and with it the number of barriers, grows with the diameter of the graph: a BFS of the USA road network in our experiments has 6,262 levels, and on such graphs the synchronization alone can cost more than the work~\cite{pasgal}.

Wait-freedom removes this dependence. It is a property of methods: a method is \emph{wait-free} if every call to it finishes in a finite number of the calling thread's own steps, however slow the other threads are~\cite{herlihy_waitfree, aomp}. Our BFS is a method that every thread calls with the same source, and that returns only once all distances are final. If this method is wait-free, no call can wait for other threads; each must be able to finish whatever work they leave undone. The search is then complete as soon as the first call returns, even if the other threads have stopped for good. Wait-freedom has mostly been studied for concurrent data structures, where wait-free implementations have long been considered too slow for practice~\cite{kogan_fastwf}. For BFS the bar is high: the blocking implementations above are heavily tuned, and a wait-free BFS method is worth using only if it is about as fast as they are when no thread is delayed.

Three problems make it hard for a BFS method to be both wait-free and fast. First, without barriers, every thread must decide for itself when a level is complete, and a thread that has fallen behind, possibly by several levels, can still write to shared state at any moment; such late writes must never undo or corrupt the work of others. Second, helping must be cheap: when no thread is delayed, threads should rarely repeat each other's work, and when one is, the others should split what it has left rather than each redo all of it. Third, direction optimization adds decisions for a whole level: all threads must expand a level in the same direction, chosen from the size of the level, which threads may see differently, and the bottom-up step needs a test of membership in the current level that remains exact for threads that are behind.

This paper presents a wait-free method for direction-optimizing BFS that solves these problems. Its contributions are:
\begin{itemize}
    \item \textbf{A wait-free method for top-down BFS} (Section~\ref{sec:wf-topdown}). The levels form a linked list, and each level holds one bucket per thread, which only that thread appends to. A thread expands its own bucket in order, publishing how far it has got, and then expands what is left of the other buckets, taking the entries in a random order, so the work of a delayed thread is taken over without waiting for it. Every write that lets other threads skip work, such as setting a distance, marking a vertex expanded or publishing progress, comes after the work it covers, so a thread that stops at any point leaves nothing that the others cannot redo.
    \item \textbf{Direction optimization that keeps the method wait-free} (Section~\ref{sec:dir-opt}). The threads agree on the direction of each level with a single compare-and-swap: each proposes a plan from what it sees, and all adopt the first proposal to claim the level. Bottom-up levels divide the vertices into blocks, which are owned and helped like buckets. A one-byte code per vertex, distinct for every level and never reused within a search, gives the bottom-up step an exact test of the current level, even for a thread that has fallen behind.
    \item \textbf{Proofs} (Section~\ref{sec:proof}) that every vertex receives its BFS distance, that the distances are final once any call returns, and that every call returns after a bounded number of the calling thread's own steps. The proofs cover both directions; the top-down method is the special case in which every level goes top-down.
    \item \textbf{A variant that bounds memory} (Section~\ref{sec:variants}), by reusing two level nodes, one for the even levels and one for the odd, instead of allocating a node per level, so that its memory depends on the largest levels rather than on all of them. Tags on the reused structures keep it correct, and wait-free, when threads fall behind.
    \item \textbf{An evaluation} (Section~\ref{sec:experiments}) on four real-world graphs and one synthetic graph, on a 128-core, two-socket Intel Xeon machine, against GAPBS, GBBS and PASGAL. Comparing each program's best time over thread counts and placements, our BFS is 1.33--1.89$\times$ faster than the fastest of the three on four of the graphs. On the road network, it is within 2\% of PASGAL, the fastest there, and the variant with two level nodes is 1.09$\times$ faster than PASGAL. Against GAPBS alone, our BFS is 1.33--2.26$\times$ faster on all five graphs. On LiveJournal, when one of 16 threads stalls for 10~ms at the start of every level or round, GAPBS, GBBS and PASGAL take 2.8--9.0$\times$ as long as without stalls, while our BFS takes 16\% longer.
\end{itemize}
To the best of our knowledge, this is the first BFS method that is wait-free for threads cooperating on a single search. Section~\ref{sec:related} discusses earlier wait-free BFS operations on concurrent graphs, each of which runs a search in a single thread, and parallel BFS methods that remove the barriers between levels without making a search wait-free.

Section~\ref{sec:background} gives background on BFS and progress guarantees and states our model. Sections~\ref{sec:wf-topdown} and~\ref{sec:dir-opt} present the method, Section~\ref{sec:proof} proves it correct and wait-free, and Section~\ref{sec:variants} describes the variant that bounds memory. Section~\ref{sec:experiments} evaluates both, Section~\ref{sec:related} discusses related work, and Section~\ref{sec:conclusion} concludes.

\section{Background and Model}
\label{sec:background}

This section reviews level-synchronous and direction-optimizing BFS, defines the progress guarantees we use, and states the model under which our method is wait-free. Throughout the paper, we say \emph{vertex} for an element of a graph and \emph{node} for an element of a linked list.

\subsection{Breadth-First Search}
\label{sec:bg-bfs}

We treat every graph as undirected. A graph $G = (V, E)$ has $n$ vertices, numbered $0$ to $n-1$, and we store each of its edges $\{u, v\}$ as the two directed edges $(u, v)$ and $(v, u)$; $m$ counts these directed edges. The graph is kept in compressed sparse row (CSR) form: the neighbors of a vertex $u$ are $\mathit{adj}[\mathit{offsets}[u]]$ up to, not including, $\mathit{adj}[\mathit{offsets}[u+1]]$, and the \emph{degree} of $u$ is $\mathit{offsets}[u+1] - \mathit{offsets}[u]$.

A BFS from a source vertex $s$ computes the distance $\delta(v)$ of every vertex $v$: the number of edges on a shortest path from $s$ to $v$, or $-1$ if there is no such path. Level $d$ is the set $L_d = \{v : \delta(v) = d\}$, so $L_0 = \{s\}$. The search keeps the distances it has found in an array $\mathit{dist}$, which holds $-1$ for every vertex at the start. A vertex is \emph{discovered} when its distance is set, and \emph{expanded} when its neighbors are examined; the level being expanded is the \emph{frontier}.

A top-down step expands the frontier $L_d$: it examines every neighbor $v$ of every vertex in $L_d$, and discovers each $v$ that is still undiscovered, setting $\mathit{dist}[v] = d+1$ and adding $v$ to $L_{d+1}$. A level-synchronous parallel BFS~\cite{bader_bfs, charles_bfs} divides each frontier among the threads, as \textsc{TopDownStep} in Algorithm~\ref{alg:do-bfs} shows. When several threads find the same undiscovered vertex, a compare-and-swap (CAS) on its distance lets only one of them add it to the next level. The step ends at a barrier, since no thread may expand $L_{d+1}$ before all of it has been discovered: otherwise a thread could reach a vertex of $L_{d+1}$ that is still undiscovered from another vertex of $L_{d+1}$, and give it the distance $d+2$.

A bottom-up step~\cite{beamer_dobfs} works the other way around: every undiscovered vertex $v$ scans its own neighbors for one in the frontier, and if it finds one, it sets $\mathit{dist}[v] = d+1$, adds $v$ to $L_{d+1}$, and stops scanning (\textsc{BottomUpStep} in Algorithm~\ref{alg:do-bfs}). Only the thread that scans $v$ writes its distance, so the step needs no CAS, but it needs a fast test of whether a vertex is in the frontier; GAPBS keeps the frontier as a bitmap for this test. When the frontier is large, a bottom-up step examines far fewer edges than a top-down step: most undiscovered vertices find a frontier neighbor among their first few neighbors, while a top-down step examines every edge of the frontier, most of which lead to vertices that are already discovered.

A direction-optimizing BFS chooses the direction of every step. The rule of Beamer et al., as GAPBS implements it, keeps two counts: $m_f$, the number of edges of the frontier, which is the sum of the degrees of its vertices, and $m_u$, an estimate of the edges not yet examined, which starts at $m$ and drops by $m_f$ with every top-down step. A search going top-down switches to bottom-up once $m_f > m_u / \alpha$. A search going bottom-up continues while the frontier grows or holds more than $n / \beta$ vertices, and switches back to top-down otherwise. GAPBS sets $\alpha = 15$ and $\beta = 18$; our method uses the same rule and values (Section~\ref{sec:dir-opt}).

\begin{algorithm}
\caption{Level-synchronous direction-optimizing BFS, as in GAPBS. Every step ends at a barrier.}
\label{alg:do-bfs}
\begin{algorithmic}[1]
\Procedure{TopDownStep}{$F$, $d$}
    \State $F' \gets \emptyset$
    \ForAll{$u \in F$ \textbf{in parallel}}
        \ForAll{neighbors $v$ of $u$}
            \If{$\mathit{dist}[v] = -1$ \textbf{and} \Call{CAS}{$\mathit{dist}[v], -1, d+1$}}
                \State add $v$ to $F'$
            \EndIf
        \EndFor
    \EndFor
    \State \textbf{barrier} \Comment{waits for the slowest thread}
    \State \Return $F'$
\EndProcedure
\Statex
\Procedure{BottomUpStep}{$F$, $d$}
    \State $F' \gets \emptyset$
    \ForAll{$v \in V$ with $\mathit{dist}[v] = -1$ \textbf{in parallel}}
        \ForAll{neighbors $u$ of $v$}
            \If{$u \in F$}
                \State $\mathit{dist}[v] \gets d+1$; add $v$ to $F'$
                \State \textbf{break}
            \EndIf
        \EndFor
    \EndFor
    \State \textbf{barrier} \Comment{waits for the slowest thread}
    \State \Return $F'$
\EndProcedure
\Statex
\Procedure{BFS}{$s$}
    \State $\mathit{dist}[s] \gets 0$; $F \gets \{s\}$; $d \gets 0$
    \While{$F \neq \emptyset$}
        \If{the direction rule chooses top-down}
            \State $F \gets$ \Call{TopDownStep}{$F$, $d$}
        \Else
            \State $F \gets$ \Call{BottomUpStep}{$F$, $d$}
        \EndIf
        \State $d \gets d+1$
    \EndWhile
\EndProcedure
\end{algorithmic}
\end{algorithm}

\subsection{Progress Guarantees}
\label{sec:bg-progress}

Progress guarantees are properties of methods~\cite{aomp}. A method is \emph{blocking} if a delay of one thread can keep the calls of other threads from finishing, as when a call waits for a lock that a delayed thread holds, or at a barrier that a delayed thread has not reached. A method is \emph{lock-free} if it guarantees that infinitely often some call finishes in a finite number of steps: the system as a whole makes progress, although a particular call may never finish. A method is \emph{wait-free} if every call finishes in a finite number of the calling thread's own steps, whatever the other threads do~\cite{herlihy_waitfree}, and \emph{bounded wait-free} if there is a bound on that number. A call of a wait-free method therefore never waits for another thread, and even a thread that stops for good delays no other. The usual way to make a method wait-free is \emph{helping}: instead of waiting for a delayed thread to finish its part of the work, a thread finishes that part itself~\cite{aomp, kogan_fastwf}.

The BFS method of Algorithm~\ref{alg:do-bfs} is blocking: a thread that is delayed before the barrier of a step delays every other thread there, and one that stops for good keeps them waiting forever.

\subsection{System Model}
\label{sec:bg-model}

There are $T$ threads, numbered $0$ to $T-1$, which communicate through shared memory. The number of threads is fixed before a search, and every thread knows it and its own number. Threads are asynchronous: any thread can be delayed for any length of time between any two of its steps, or stop for good, and we assume nothing about their relative speeds. The graph does not change during a search.

Memory accesses follow total store order (TSO), the memory model of x86 processors~\cite{x86tso, intelmanual, amdmanual}. Aligned loads and stores of up to eight bytes are atomic, and so is a compare-and-swap (CAS) on such a location. The stores of a thread become visible to all other threads at the same time, and in the order in which the thread made them, and the loads of a thread take effect in program order. The one reordering that TSO allows is that a load can take effect before an earlier store of the same thread to another location becomes visible, while the store still waits in the thread's store buffer. Our method relies only on the order between two stores of a thread and between two loads of a thread, and on CAS, which x86 performs as a locked instruction that also orders the accesses around it. We write the method with ordinary loads and stores, and put compiler barriers where their order matters, so that the compiler keeps those accesses in program order. On a weaker memory model, the stores that must stay ordered would need release semantics, and the loads that read them acquire semantics.

Each thread allocates the memory it needs during a search from its own arena, which hands out blocks by moving a pointer and asks the system for more memory only when it runs out. We assume that a thread can allocate memory in a bounded number of its own steps.

\subsection{The BFS Method}
\label{sec:bg-method}

Every thread calls the method \textsc{BFS}$(s)$ once per search, with the same source $s$. Before the first call, $\mathit{dist}[v] = -1$ for every vertex $v$, and the method's other per-vertex arrays hold their initial values. Resetting these arrays between searches is not part of the method; since a search reuses them, the next search starts only once every call of the previous one has returned.

When any call returns, $\mathit{dist}[v] = \delta(v)$ for every vertex $v$, and no step that any thread takes afterwards changes $\mathit{dist}$. A search is therefore complete when its first call returns, and that is when we stop timing it (Section~\ref{sec:experiments}). Section~\ref{sec:proof} proves these properties, and that every call returns after a bounded number of the calling thread's own steps.

\section{A Wait-Free Method for Top-Down BFS}
\label{sec:wf-topdown}

This section presents our BFS method with top-down steps only; Section~\ref{sec:dir-opt} adds bottom-up steps. The method replaces the barrier at the end of each level with bookkeeping that lets every thread decide on its own when a level is complete, and finish any part of it that other threads have left undone. Algorithm~\ref{alg:wf-td} gives the method as thread $t$ runs it.

\subsection{Levels and Buckets}
\label{sec:td-structure}

The levels of a search form a linked list, whose node $N_d$ describes level $d$ (Figure~\ref{fig:td-example}). A node holds its depth $d$, a pointer $\mathit{next}$ to the node of level $d+1$, which is null until some thread links that node, and one \emph{entry} for each thread. Entry $t$ of $N_d$ holds the \emph{bucket} $B_{d,t}$, a list of vertices of level $d$ that thread $t$ has discovered, and two fields that record how far the expansion of the bucket has got: $\mathit{last}_{d,t}$, the position of the last vertex of $B_{d,t}$ that thread $t$ itself has finished, which starts at $-1$, and the flag $\mathit{done}_{d,t}$, which a thread sets once it has finished the vertices of the bucket. Only thread $t$ appends to $B_{d,t}$, and only thread $t$ writes $\mathit{last}_{d,t}$.

Besides the distances, the method keeps a flag $\mathit{vis}[u]$ for every vertex $u$, which marks $u$ as expanded: a thread sets it once it has examined every neighbor of $u$, and other threads then skip $u$. A vertex can be in several buckets of a level, since several threads can discover it at the same time, but every copy is in the right level: a vertex in $B_{d,t}$ is at distance $d$. Before a search, $N_0$ holds the source $s$ in the bucket of thread $0$, $\mathit{dist}[s] = 0$, and every other distance is $-1$ and every flag unset.

\begin{algorithm}
\caption{Wait-free top-down BFS, as thread $t$ runs it. For a node $N = N_d$, $N.B[o]$, $N.\mathit{last}[o]$ and $N.\mathit{done}[o]$ stand for $B_{d,o}$, $\mathit{last}_{d,o}$ and $\mathit{done}_{d,o}$.}
\label{alg:wf-td}
\begin{algorithmic}[1]
\Procedure{BFS}{$s$}
    \State $N \gets N_0$
    \While{$N \neq \mathit{null}$}
        \State \textbf{fence} \label{ln:td-fence}
        \State \Call{ExpandLevel}{$N$}
        \State $N \gets N.\mathit{next}$ \label{ln:td-next}
    \EndWhile
\EndProcedure
\Statex
\Procedure{ExpandLevel}{$N$}
    \For{$k \gets 0$ \textbf{to} $T-1$}
        \State $o \gets (t + k) \bmod T$
        \IfThen{$N.\mathit{done}[o]$ \textbf{or} $N.B[o]$ is empty}{\textbf{continue}}
        \IfThen{$N.\mathit{next} = \mathit{null}$}{\Call{AddLevel}{$N$}} \label{ln:td-add}
        \State $\mathit{size} \gets$ the size of $N.B[o]$ \label{ln:td-size}
        \If{$o = t$}
            \For{$i \gets 0$ \textbf{to} $\mathit{size}-1$} \label{ln:td-own}
                \State \Call{Expand}{$N$, $N.B[t][i]$}
                \State $N.\mathit{last}[t] \gets i$ \label{ln:td-last}
            \EndFor
        \Else
            \State $\mathit{first} \gets N.\mathit{last}[o] + 1$ \label{ln:td-first}
            \ForAll{$i$ from $\mathit{first}$ to $\mathit{size}-1$, in a random order} \label{ln:td-help}
                \State \Call{Expand}{$N$, $N.B[o][i]$}
            \EndFor
        \EndIf
        \State $N.\mathit{done}[o] \gets \mathit{true}$ \label{ln:td-done}
    \EndFor
\EndProcedure
\Statex
\Procedure{Expand}{$N$, $u$}
    \IfThen{$\mathit{vis}[u]$}{\textbf{return}}
    \ForAll{neighbors $v$ of $u$}
        \If{$\mathit{dist}[v] = -1$}
            \State append $v$ to $N.\mathit{next}.B[t]$ \label{ln:td-push}
            \State $\mathit{dist}[v] \gets N.\mathit{depth} + 1$ \label{ln:td-dist}
        \EndIf
    \EndFor
    \State $\mathit{vis}[u] \gets \mathit{true}$ \label{ln:td-vis}
\EndProcedure
\Statex
\Procedure{AddLevel}{$N$}
    \State $M \gets$ the spare node of $t$, or a new node
    \State $M.\mathit{depth} \gets N.\mathit{depth} + 1$
    \If{\Call{CAS}{$N.\mathit{next}$, $\mathit{null}$, $M$}} \label{ln:td-cas}
        \State $t$ has no spare node
    \Else
        \State $M$ becomes the spare node of $t$
    \EndIf
\EndProcedure
\end{algorithmic}
\end{algorithm}

\subsection{Expanding a Level}
\label{sec:td-expand}

A thread expands level $d$ by visiting the $T$ entries of $N_d$, starting with its own and continuing with those of threads $t+1$, $t+2$, and so on, modulo $T$ (\textsc{ExpandLevel}). It skips an entry whose bucket is empty or whose flag $\mathit{done}$ is set. Otherwise, it reads the size of the bucket, expands the vertices before that position, sets $\mathit{done}$, and moves on to the next entry.

In its own bucket, a thread expands the vertices in order, and after each one it publishes its position in $\mathit{last}_{d,t}$ (line~\ref{ln:td-last}). In the bucket of another thread $o$, it expands only the vertices after position $\mathit{last}_{d,o}$, the ones $o$ had not finished when the helper read it, and takes them in a random order (line~\ref{ln:td-help}). Expanding a vertex $u$ does nothing if $\mathit{vis}[u]$ is set (\textsc{Expand}). Otherwise, the thread examines every neighbor $v$ of $u$, and if $\mathit{dist}[v] = -1$, it appends $v$ to its own bucket in $N_{d+1}$ and sets $\mathit{dist}[v] = d+1$. After the last neighbor, it sets $\mathit{vis}[u]$.

A thread thus works on its own bucket first, and helps the others once it is done with it. When no thread is delayed, the threads mostly expand their own buckets, and a thread that finishes early takes over part of the work of those still busy. A thread that is delayed, or has stopped for good, holds no one up: the others expand what is left of its bucket, set its flag $\mathit{done}$, and move on.

When a thread has visited all $T$ entries of $N_d$, every vertex of level $d$ has been expanded, by it or by another thread, so every vertex of level $d+1$ has been discovered and is in some bucket of $N_{d+1}$. The thread then moves on to $N_{d+1}$ without waiting for anyone (line~\ref{ln:td-next}). Before it reads the buckets of the new level, it executes a fence (line~\ref{ln:td-fence}), which makes all of its earlier writes visible to the other threads. On x86, its last appends to $N_{d+1}$ could otherwise still sit in its store buffer, where only it sees them, while other threads read the buckets of $N_{d+1}$ (Section~\ref{sec:pf-tso}); the fence costs one instruction per level. In the method of Section~\ref{sec:dir-opt}, the CAS that claims the plan of a level, which comes before any thread expands the level, does the same job, so that method needs no fence. Slower threads may still be expanding level $d$; Section~\ref{sec:td-order} explains why they cannot harm the levels after it, and Section~\ref{sec:proof} proves it.

\subsection{The Order of Writes}
\label{sec:td-order}

Threads skip work on the strength of what other threads have written. A vertex whose distance is set is not discovered again, a vertex whose flag $\mathit{vis}$ is set is not expanded again, helpers do not expand the vertices up to position $\mathit{last}_{d,o}$ of a bucket, and no thread visits a bucket whose flag $\mathit{done}$ is set. Each of these writes comes after the work it lets others skip, so a thread that stops between any two of its steps leaves no work undone that others would skip:
\begin{itemize}
    \item A thread appends $v$ to its bucket, and only then sets $\mathit{dist}[v]$ (lines~\ref{ln:td-push} and~\ref{ln:td-dist}). If it stops in between, $\mathit{dist}[v]$ is still $-1$, and another thread discovers $v$ again and appends it to its own bucket. In the other order, a thread that stopped after setting $\mathit{dist}[v]$ would leave $v$ discovered but in no bucket, and no thread would ever expand it.
    \item A thread sets $\mathit{vis}[u]$ only after it has examined every neighbor of $u$ (line~\ref{ln:td-vis}). A thread that stops in the middle of expanding $u$ leaves $\mathit{vis}[u]$ unset, so other threads expand $u$ again, from its first neighbor.
    \item A thread publishes $\mathit{last}_{d,t} = i$ only after it has finished the vertex at position $i$ of its bucket, and sets $\mathit{done}$ only after it has finished every vertex it read in the bucket (lines~\ref{ln:td-last} and~\ref{ln:td-done}).
\end{itemize}
Every thread that discovers $v$ writes the same distance, $d+1$, so a plain store suffices for it. Section~\ref{sec:dir-opt} replaces the store with a CAS, so that the direction rule counts the edges of each vertex only once.

A delayed thread can also resume long after the others have moved on. Suppose thread $t$ stopped while expanding a vertex $u$ of level $d$, after it read $\mathit{dist}[v] = -1$ for a neighbor $v$ but before it appended $v$. Before any thread could finish level $d$, some thread finished expanding $u$, and in doing so either appended $v$ itself or found it already discovered, and therefore already appended, by another thread. When thread $t$ resumes, it appends $v$ to $B_{d+1,t}$, which the others may have expanded already, and writes $\mathit{dist}[v] = d+1$ again. Neither write does harm: the distance is the one already there, and the late copy of $v$ is in the right level and duplicates one that the others have expanded or will expand, so a thread that read $B_{d+1,t}$ before the late append missed nothing. Such late appends are why a vertex can have several copies in a level, and why $\mathit{done}$ means that the vertices a thread read in a bucket are finished, rather than every vertex the bucket will ever hold.

\subsection{Buckets}
\label{sec:td-buckets}

A bucket has to let its owner append vertices while any number of other threads read it, without locks, and without ever letting a reader see a position that has not been written or memory that has been freed. The readers are helpers, which can read a bucket while its owner is still appending to it, and because of late appends, an owner can append to a bucket long after other threads have read it and moved on. The number of vertices a bucket will hold is not known in advance, and helpers read a bucket by position, starting after $\mathit{last}$ and taking blocks in a random order. A bucket is therefore a dynamic array, and since only its owner writes it, appends need no CAS.

\begin{algorithm}
\caption{Appending to a bucket $B$, which only its owner does, and reading it, which any thread can do.}
\label{alg:bucket}
\begin{algorithmic}[1]
\Procedure{Append}{$B$, $v$}
    \If{$B.\mathit{size} = B.\mathit{capacity}$}
        \State $c \gets \max(2 \cdot B.\mathit{capacity}, c_0)$
        \State $A \gets$ an array of $c$ vertices from the owner's arena
        \State copy $B.\mathit{array}[0], \ldots, B.\mathit{array}[B.\mathit{size}-1]$ into $A$
        \State $B.\mathit{array} \gets A$ \label{ln:bk-array}
        \State $B.\mathit{capacity} \gets c$
    \EndIf
    \State $B.\mathit{array}[B.\mathit{size}] \gets v$
    \State $B.\mathit{size} \gets B.\mathit{size} + 1$ \label{ln:bk-size}
\EndProcedure
\Statex
\Procedure{Read}{$B$}
    \State $s \gets B.\mathit{size}$
    \State $A \gets B.\mathit{array}$
    \State \Return $A[0], \ldots, A[s-1]$
\EndProcedure
\end{algorithmic}
\end{algorithm}

A bucket has three fields: the pointer $\mathit{array}$ to its current array, the number $\mathit{size}$ of vertices appended so far, and the $\mathit{capacity}$ of the current array, which only the owner reads (Algorithm~\ref{alg:bucket}). When the array is full, the owner allocates one twice as large from its arena, copies the vertices into it, and only then publishes it in $\mathit{array}$ (line~\ref{ln:bk-array}); the first array has a small fixed size $c_0$. To append a vertex, the owner writes it at position $\mathit{size}$ of the array, and only then increases $\mathit{size}$ (line~\ref{ln:bk-size}). A reader reads in the opposite order: $\mathit{size}$ first, then $\mathit{array}$, and then the vertices before position $\mathit{size}$. In Algorithm~\ref{alg:wf-td}, a thread reads the size of a bucket at line~\ref{ln:td-size}, and its array only afterwards.

These two orders keep every read within positions that have been written, even when the read overlaps an append that grows the array. Suppose a reader reads the size $s$. Since the owner publishes an array before any size that needs it, and the reader reads the size before the array, the array the reader then finds is the one that was current when $s$ was published, or a later one. The first holds the first $s$ vertices, which are never moved or overwritten, and every later array holds copies of them, made before it was published. Whichever array the reader finds, its first $s$ positions are written and hold the same vertices. A reader can miss vertices appended after it read the size, but that is of no consequence: by the time a thread reads the buckets of $N_d$, every vertex of level $d$ is in them already, and the appends that come later are late copies (Section~\ref{sec:td-order}).

An array that the owner has replaced cannot be freed during the search, since a reader may still hold its pointer, and a delayed reader may hold it for any length of time. Finding out when no reader holds it any more would take extra work from every reader, so old arrays instead stay in the owner's arena until the search is over, and the arenas are reset as a whole once every call has returned. Since capacities double, the old arrays of a bucket together hold fewer positions than its current array, and a bucket of $k$ vertices that has outgrown its first array takes fewer than $4k$ positions in all. An append takes constant time amortized, and the copy in a growth moves at most $n$ vertices, since a thread appends each vertex at most once in a search (Section~\ref{sec:proof}).

\subsection{Linking Levels and Ending the Search}
\label{sec:td-link}

A thread that finds a bucket of $N_d$ to expand first makes sure that $N_{d+1}$ exists, since the vertices it discovers go there (line~\ref{ln:td-add}). If $N_d.\mathit{next}$ is null, the thread prepares a node for level $d+1$ and tries to link it with a CAS on $N_d.\mathit{next}$ from null (line~\ref{ln:td-cas}). If the CAS fails, another thread has linked a node already, and every thread uses that one. A node that lost the race was never seen by any other thread, so its thread keeps it as a spare and tries it again at the next level, instead of allocating another. Threads tend to reach a level together, and nearly all of them prepare a node for the next one, so without spares most of these nodes would be wasted.

A thread links $N_{d+1}$ before it expands any vertex of level $d$, and sets a flag $\mathit{done}$ only afterwards, so a thread that skips a done bucket still finds $N_{d+1}$ linked. A level without vertices, on the other hand, gets no successor: all of its buckets are empty, so no thread links a node after it, and every thread that visits it finds $\mathit{next}$ null and returns (line~\ref{ln:td-next}). The first such level is level $D+1$, where $D$ is the largest distance from $s$, so every call of the method goes through the nodes of levels $0$ to $D+1$ and returns there.

\subsection{An Example}
\label{sec:td-example}

\begin{figure}
    \centering
    \begin{tikzpicture}[vertex/.style={circle, draw, minimum size=0.55cm, inner sep=0pt}]
        \node[vertex] (A) at (0,0) {$A$};
        \node[vertex] (B) at (1.5,0.9) {$B$};
        \node[vertex] (C) at (1.5,0) {$C$};
        \node[vertex] (D) at (1.5,-0.9) {$D$};
        \node[vertex] (E) at (3,0.45) {$E$};
        \node[vertex] (F) at (3,-0.45) {$F$};
        \draw (A) -- (B);
        \draw (A) -- (C);
        \draw (A) -- (D);
        \draw (B) -- (E);
        \draw (C) -- (E);
        \draw (C) -- (F);
    \end{tikzpicture}

    \vspace{0.3cm}

    \begin{tikzpicture}[>=Stealth, levelnode/.style={draw, rectangle, align=left, font=\footnotesize, minimum width=1.55cm, inner sep=3pt}]
        \node[levelnode] (n0) at (0,0) {$N_0$\\ $0$: $A$\\ $1$: --};
        \node[levelnode] (n1) at (2.05,0) {$N_1$\\ $0$: $B, C$\\ $1$: $C, D$};
        \node[levelnode] (n2) at (4.1,0) {$N_2$\\ $0$: $E$, \textcolor{gray}{$F$}\\ $1$: $F$};
        \node[levelnode] (n3) at (6.15,0) {$N_3$\\ $0$: --\\ $1$: --};
        \draw[->] (n0) -- (n1);
        \draw[->] (n1) -- (n2);
        \draw[->] (n2) -- (n3);
    \end{tikzpicture}
    \caption{A search from $A$ with two threads. Each node lists the buckets of threads $0$ and $1$. Both threads discover $C$, so level 1 holds it twice. Thread $0$ stalls while expanding $C$, and appends $F$ to its bucket in $N_2$ (in gray) only after thread $1$ has finished the search.}
    \Description{A graph with vertices A to F, and the four level nodes of a search from A by two threads, with the buckets of each thread in each node.}
    \label{fig:td-example}
\end{figure}

Figure~\ref{fig:td-example} shows a search from $A$ with two threads. Level 0 holds $A$ in the bucket of thread 0. Thread 1 finds its own bucket empty and helps thread 0, which has not finished $A$ yet, so both threads expand $A$ at the same time. Thread 0 discovers $B$ and $C$, and thread 1, which reads $\mathit{dist}[C]$ just before thread 0 sets it, discovers $C$ and $D$, so level 1 holds $C$ twice. In level 1, thread 0 expands $B$, which discovers $E$, and publishes $\mathit{last}_{1,0} = 0$. It then starts on $C$, reads $\mathit{dist}[F] = -1$, and stalls before it appends $F$. Thread 1 meanwhile expands its own bucket: $C$, which is not marked yet, so thread 1 discovers $F$, and then $D$, which discovers nothing new. It then helps thread 0 with the vertices after position 0 of $B_{1,0}$. The only one is $C$, which thread 1 has just expanded, so it skips it and sets $\mathit{done}_{1,0}$. Thread 1 goes on to expand level 2, which links $N_3$, finds the buckets of $N_3$ empty, and returns. The search is now complete, although thread 0 is still stalled in level 1. When thread 0 resumes, it appends $F$ to $B_{2,0}$, writes $\mathit{dist}[F] = 2$ again, and finishes $C$. It then finds the buckets of $N_2$ done and those of $N_3$ empty, and returns too.

\subsection{Keeping Helpers Apart}
\label{sec:td-contention}

The method is correct whatever order the threads take entries and vertices in; the orders we use keep the threads out of each other's way. Owners walk their buckets in order, which keeps their reads sequential. A helper starts after the owner's published position, and walks the rest of the bucket in small blocks of consecutive entries, taking the blocks in a random order of its own and the vertices of a block in order. Helpers on the same bucket thus rarely expand the same vertex at the same time, or the vertex its owner is on, while each still reads consecutive entries. A vertex with many neighbors raises the same problem within a single expansion: if the owner and helpers expanded it together, walking its neighbors in the same order, they would all write the same entries of $\mathit{dist}$ at the same moment. A helper therefore walks a long neighbor list in small blocks as well, in a random order of its own. Each random order is an affine permutation of the $k$ blocks, which visits block $(a + i b) \bmod k$ in step $i$, for a random $a$ and a random $b$ coprime to $k$, so it takes two random numbers and no memory.

\section{Direction Optimization That Keeps the Method Wait-Free}
\label{sec:dir-opt}

Adding bottom-up levels to the method of Section~\ref{sec:wf-topdown} raises three problems that a blocking implementation solves at its barriers. First, the direction is a decision for a whole level, which every thread must make the same way, although threads see the level at different times and can count it differently. Second, a bottom-up level works over all undiscovered vertices rather than over buckets, so it needs its own way of dividing the work and of taking over the work of a delayed thread. Third, a bottom-up level tests whether a vertex is in the frontier, and the test must stay exact for a thread that has fallen behind and runs it after the others have moved on. This section takes the three in turn. Algorithm~\ref{alg:wf-do} shows how a thread chooses the direction of each level, and Algorithm~\ref{alg:wf-bu} how it expands a bottom-up level; top-down levels are expanded as in Algorithm~\ref{alg:wf-td}, except for one change in how vertices are discovered.

\begin{algorithm}
\caption{Direction-optimizing BFS, as thread $t$ runs it. \textsc{ExpandLevel} and \textsc{AddLevel} are those of Algorithm~\ref{alg:wf-td}, with \textsc{Discover} in place of lines~\ref{ln:td-push} and~\ref{ln:td-dist}.}
\label{alg:wf-do}
\begin{algorithmic}[1]
\Procedure{BFS}{$s$}
    \State $N \gets N_0$; $P \gets$ the plan before level $0$
    \While{$N \neq \mathit{null}$}
        \State $P' \gets$ \Call{ChoosePlan}{$N$, $P$}
        \If{$P'.\mathit{direction}$ is top-down}
            \State \Call{ExpandLevel}{$N$}
        \Else
            \State \Call{ScanLevel}{$N$, $P'$, $P$}
        \EndIf
        \State $P \gets P'$; $N \gets N.\mathit{next}$
    \EndWhile
\EndProcedure
\Statex
\Procedure{ChoosePlan}{$N$, $P$}
    \If{$N.\mathit{planner}$ is empty}
        \State $N.\mathit{proposal}[t] \gets$ \Call{Rule}{$P$, $N$} \label{ln:do-propose}
        \State \Call{CAS}{$N.\mathit{planner}$, empty, $t$} \label{ln:do-claim}
    \EndIf
    \State \Return $N.\mathit{proposal}[N.\mathit{planner}]$ \label{ln:do-adopt}
\EndProcedure
\Statex
\Procedure{Discover}{$N$, $v$} \Comment{in top-down levels}
    \State append $v$ to $N.\mathit{next}.B[t]$
    \If{\Call{CAS}{$\mathit{dist}[v]$, $-1$, $N.\mathit{depth}+1$}} \label{ln:do-cas}
        \State add the degree of $v$ to $N.\mathit{next}.\mathit{edges}[t]$
    \EndIf
\EndProcedure
\end{algorithmic}
\end{algorithm}

\subsection{Agreeing on a Direction}
\label{sec:do-plan}

Before it expands a level, a thread must know the level's \emph{plan}: its direction, and for a bottom-up level the code of its frontier (Section~\ref{sec:do-codes}), together with the counts that the direction rule of Section~\ref{sec:bg-bfs} carries from one level to the next, such as $m_u$ and the size of the previous frontier. All threads must use the same plan for a level. In a top-down level, the fields $\mathit{last}$ and $\mathit{done}$ of an entry refer to positions in its bucket, while in a bottom-up level they refer to blocks of vertices, so a thread that scanned blocks while another expanded buckets would misread the other's progress, and could skip work that no thread has done. A bottom-up level must also test the frontier with the code that was written for it.

A blocking implementation lets one thread decide and makes the others wait for the decision. Instead, each node $N_d$ has a field $\mathit{planner}$, which starts out empty, and each of its entries has room for a \emph{proposal} (\textsc{ChoosePlan}). A thread that enters $N_d$ and finds $\mathit{planner}$ empty computes a plan with the direction rule (\textsc{Rule}) from what it sees: the frontier size $n_f$, which is the sum of the sizes of the buckets of $N_d$; the count $m_f$, which is the sum of the edge counts of its entries, described below; and the plan of level $d-1$. It writes the plan into its own entry (line~\ref{ln:do-propose}) and then tries to claim the level with a CAS on $\mathit{planner}$ from empty to its own number (line~\ref{ln:do-claim}). Whether its CAS succeeds or not, it then adopts the proposal of the thread that $\mathit{planner}$ names, as does every thread that finds $\mathit{planner}$ already set (line~\ref{ln:do-adopt}). A thread makes at most one CAS per level, and none waits for another: a thread claims a level only after it has written its proposal, so the proposal is there for every thread that reads the claim.

Proposals can differ, since threads read the buckets at different times and late appends change their sizes, but the CAS makes every thread adopt the same one. Which one wins does not matter for correctness, since the method finds the same distances in either direction; the rule only affects speed, and proposals made from the same buckets differ only by the few late copies. Since all threads adopt the same plan for every level, and start from the same plan before level 0, they all propose from the same plan of the level before.

The rule needs $m_f$, the number of edges of the frontier, and it gets it from the top-down level that discovered the frontier. In the method of Section~\ref{sec:wf-topdown}, every thread that finds a vertex undiscovered appends it and writes its distance, so a vertex discovered by two threads at once would have its edges counted twice. Counted that way, $m_f$ would overstate the edges of a large frontier, which can make the rule switch to bottom-up a level too early, and $m_u$, which drops by $m_f$ with every top-down level, would run down too fast. Top-down levels therefore discover a vertex with a CAS on its distance from $-1$ to $d+1$, still after appending it (\textsc{Discover}), and only the thread whose CAS succeeds adds the degree of the vertex to the edge count of its entry in $N_{d+1}$. A thread whose CAS fails has found the vertex discovered by another thread, which appended it first, so the order of Section~\ref{sec:td-order} still holds. After a bottom-up level, the rule takes $m_f = 1$, as GAPBS does.

\begin{algorithm}
\caption{A bottom-up level, as thread $t$ runs it. $P'$ is the plan of the level and $P$ that of the level before; $X_{o,k}$ is the $k$-th block of thread $o$, and $K_o$ the number of its blocks.}
\label{alg:wf-bu}
\begin{algorithmic}[1]
\Procedure{ScanLevel}{$N$, $P'$, $P$}
    \IfThen{$N.\mathit{next} = \mathit{null}$}{\Call{AddLevel}{$N$}}
    \If{$P'.\mathit{code} > 0$ \textbf{and} $P.\mathit{direction}$ is top-down}
        \State \Call{WriteCodes}{$N$, $P'.\mathit{code}$} \label{ln:bu-codes}
    \EndIf
    \For{$k \gets 0$ \textbf{to} $K_t - 1$, \textbf{while not} $N.\mathit{done}[t]$} \label{ln:bu-own}
        \IfThen{$\mathit{stamp}[t][k] \neq N.\mathit{depth}$}{\Call{ScanBlock}{$N$, $P'$, $X_{t,k}$}}
        \State $N.\mathit{last}[t] \gets k$ \label{ln:bu-last}
    \EndFor
    \State $N.\mathit{done}[t] \gets \mathit{true}$
    \For{$j \gets 1$ \textbf{to} $T-1$}
        \State $o \gets (t + j) \bmod T$
        \ForAll{blocks $k > N.\mathit{last}[o]$, in a random order}\label{ln:bu-help}
            \IfThen{$N.\mathit{done}[o]$}{\textbf{break}}
            \If{$k > N.\mathit{last}[o]$ \textbf{and} $\mathit{stamp}[o][k] \neq N.\mathit{depth}$}
                \State \Call{ScanBlock}{$N$, $P'$, $X_{o,k}$}
                \State $\mathit{stamp}[o][k] \gets N.\mathit{depth}$ \label{ln:bu-stamp}
            \EndIf
        \EndFor
        \State $N.\mathit{done}[o] \gets \mathit{true}$
    \EndFor
\EndProcedure
\Statex
\Procedure{ScanBlock}{$N$, $P'$, $X$}
    \ForAll{$v \in X$ with $\mathit{dist}[v] = -1$}
        \If{some neighbor $u$ of $v$ has $\mathit{code}[u] = P'.\mathit{code}$} \label{ln:bu-test}
            \State append $v$ to $N.\mathit{next}.B[t]$ \label{ln:bu-append}
            \State $\mathit{code}[v] \gets P'.\mathit{code} + 1$ \label{ln:bu-code}
            \State $\mathit{dist}[v] \gets N.\mathit{depth} + 1$ \label{ln:bu-dist}
        \EndIf
    \EndFor
\EndProcedure
\Statex
\Procedure{WriteCodes}{$N$, $c$}
    \For{$k \gets 0$ \textbf{to} $T-1$}
        \State $o \gets (t + k) \bmod T$
        \IfThen{$N.\mathit{codesDone}[o]$}{\textbf{continue}}
        \State $\mathit{size} \gets$ the size of $N.B[o]$
        \If{$o = t$}
            \For{$i \gets 0$ \textbf{to} $\mathit{size}-1$}
                \State $\mathit{code}[N.B[t][i]] \gets c$; $N.\mathit{codesLast}[t] \gets i$
            \EndFor
        \Else
            \For{$i \gets N.\mathit{codesLast}[o]+1$ \textbf{to} $\mathit{size}-1$}
                \State $\mathit{code}[N.B[o][i]] \gets c$
            \EndFor
        \EndIf
        \State $N.\mathit{codesDone}[o] \gets \mathit{true}$
    \EndFor
\EndProcedure
\end{algorithmic}
\end{algorithm}

\subsection{Bottom-Up Levels}
\label{sec:do-bu}

A bottom-up level divides the vertices into blocks of consecutive vertices, all of the same size, and gives block $b$ to thread $b \bmod T$, so that the blocks of every thread spread over the whole range of vertex numbers, and a level whose undiscovered vertices lie mostly in one part of the range still gives every thread a share. Blocks take the place of buckets (\textsc{ScanLevel}). A thread scans its own blocks in order, publishing in $\mathit{last}_{d,t}$ the number of the last one it has finished (line~\ref{ln:bu-last}), and then helps the other threads with the blocks after their $\mathit{last}$, in a random order of its own (line~\ref{ln:bu-help}). The flag $\mathit{done}_{d,o}$ is set once all of the blocks of thread $o$ are finished, by $o$ after its last block, or by a helper that has got through the rest; the owner and its helpers stop as soon as they see it set.

A helper cannot rely on the owner's $\mathit{last}$ alone, since a delayed owner stops moving it, and every helper of a delayed owner would then scan all of the owner's remaining blocks. A helper therefore stamps every block it finishes with the depth of the level (line~\ref{ln:bu-stamp}), and the owner and the other helpers skip a block whose stamp equals the current depth. Stamps outlive levels: a stamp left by an earlier level holds an older depth, and a thread that has fallen behind can only write the depth of its own level, which is older than the current one. Either way, a stale stamp makes a thread scan a block again, but never skip one that needs scanning.

Scanning a block means testing every vertex $v$ of the block that is still undiscovered (\textsc{ScanBlock}). If $v$ has a neighbor in the frontier, it is at distance $d+1$: the thread appends $v$ to its own bucket in $N_{d+1}$, writes the code of level $d+1$ for $v$, and only then sets $\mathit{dist}[v] = d+1$ (lines~\ref{ln:bu-append} to~\ref{ln:bu-dist}). A thread that finds $v$ discovered therefore finds it in a bucket and with its code. Threads that scan a block at the same time can discover the same vertex, but they all write the same values, so plain stores suffice. Since the discoveries of a bottom-up level go into buckets, as those of a top-down level do, the next level can go in either direction, and its plan can count its frontier.

\subsection{Frontier Codes}
\label{sec:do-codes}

The bottom-up step tests, for neighbor after neighbor, whether a vertex $u$ is in the frontier. The test could read $\mathit{dist}[u] = d$, but it is the step's most frequent memory access, and $\mathit{dist}$ has four bytes per vertex. GAPBS keeps the frontier in a bitmap instead, which it rebuilds for every bottom-up step. A bitmap reused from one level to the next would not do for our method: a thread that has fallen behind would test the frontier of level $d$ against bits that already describe a later level. We keep one byte per vertex, $\mathit{code}[v]$, and give each level that a bottom-up level tests a code of its own, $c_d$, which the plan assigns. Codes are handed out in increasing order and never reused within a search, and a vertex only ever receives the code of its own level. The test $\mathit{code}[u] = c_d$ (line~\ref{ln:bu-test}) is therefore exact at any time, for any thread, however far it has fallen behind.

A bottom-up level $d$ gets the codes of its frontier in one of two ways. If level $d-1$ was bottom-up too, its discoveries wrote $c_d$, which the plan of level $d-1$ reserved, before their distances, so every vertex of level $d$ has its code once level $d-1$ is complete. If level $d-1$ was top-down, whose appends write no codes, then every thread starts level $d$ by writing $c_d$ for the vertices in the buckets of $N_d$ (\textsc{WriteCodes}, line~\ref{ln:bu-codes}). This pass is shared the way an expansion is: a thread writes the codes of its own bucket in order, publishing how far it has got, then writes the rest of every other bucket whose codes are not done yet, and marks a bucket done once its codes are all written. Each entry has a separate pair of fields for this pass, $\mathit{codesLast}$ and $\mathit{codesDone}$, so that it does not disturb the progress of the expansion. A thread starts testing only after the pass, when every vertex of level $d$ has its code. Codes written twice or late are still right, since every vertex in the buckets of $N_d$ is in level $d$. Top-down levels write no codes, since many threads appending nearby vertices would contend for the same cache lines of $\mathit{code}$.

Codes fit in a byte, so once the plans have handed out code 255, later bottom-up levels get the code 0, which is the code of a vertex without one, and test $\mathit{dist}[u] = d$ instead; the discoveries of a level with code 255 write the code 0.

\subsection{Ending the Search}
\label{sec:do-end}

As in Section~\ref{sec:wf-topdown}, the search ends at the first level without vertices, and only a top-down level ends it. A bottom-up level always links a successor, since it cannot tell in advance whether it will discover anything, and the plan of a level whose buckets are all empty is always top-down, since every proposal for it finds $n_f = 0$. Every thread that enters that level finds its buckets empty, links no successor, and returns.

\section{Correctness and Wait-Freedom}
\label{sec:proof}

This section proves that the method of Section~\ref{sec:dir-opt} meets the specification of Section~\ref{sec:bg-method}: when any call of \textsc{BFS}$(s)$ returns, $\mathit{dist}$ holds the distances from $s$, and no later step changes them (Theorem~\ref{thm:correct}), and every call returns after a bounded number of the calling thread's own steps (Theorem~\ref{thm:waitfree}). The method is that of Algorithms~\ref{alg:wf-do} and~\ref{alg:wf-bu}, which expand top-down levels with Algorithm~\ref{alg:wf-td} and \textsc{Discover}, and keep their buckets as in Algorithm~\ref{alg:bucket}. We reason about the loads, stores, CASes and fences of the algorithms in the model of Section~\ref{sec:bg-model}, in which every thread makes them in the order that the algorithms list them; Section~\ref{sec:impl} describes how our implementation realizes this model. Sections~\ref{sec:pf-defs} to~\ref{sec:pf-waitfree} assume sequential consistency: the memory accesses of all threads take effect one at a time, in an order that keeps the program order of every thread. Section~\ref{sec:pf-tso} then shows that the proof holds under TSO as well. The proof never assumes that a thread takes another step, so it holds whatever delays the threads suffer, even if some of them stop for good.

\subsection{Definitions}
\label{sec:pf-defs}

In this section we write $\delta(v) = \infty$ for a vertex that $s$ cannot reach, so that the levels $L_0, \ldots, L_D$, where $D$ is the largest finite distance, hold exactly the vertices that $s$ reaches. Lemma~\ref{lem:list} shows that the nodes form a single list, in which $N_d$ is the node of depth $d$. A thread \emph{enters} $N_d$ when it calls \textsc{ChoosePlan}$(N_d)$, and \emph{finishes} $N_d$ when its call of \textsc{ExpandLevel}$(N_d)$ or \textsc{ScanLevel}$(N_d)$ returns; from the time its call of \textsc{ChoosePlan}$(N_d)$ returns until it finishes $N_d$, the thread is \emph{in level $d$}. Apart from its proposal, which it writes in \textsc{ChoosePlan}$(N_d)$, a thread appends to the buckets of $N_{d+1}$, writes $\mathit{dist}$, $\mathit{code}$ and $\mathit{vis}$, writes stamps, and writes the fields of the entries of $N_d$ only while it is in some level $d$, possibly long after other threads have moved on.

Let $\tau_d$ be the time of the first CAS on $N_d.\mathit{planner}$ (line~\ref{ln:do-claim}). If some thread enters $N_d$, this CAS exists, and it comes before every call of \textsc{ChoosePlan}$(N_d)$ returns, since a call returns after its own CAS, or after reading $\mathit{planner}$ set by an earlier one. Every read that a thread in level $d$ makes of the buckets of $N_d$ thus comes after $\tau_d$. The \emph{reference prefix} $S_{d,t}$ is the set of vertices at positions $0$ to $k-1$ of $B_{d,t}$, where $k$ is the size of $B_{d,t}$ at time $\tau_d$. By Lemma~\ref{lem:buckets}, these positions keep their vertices for the rest of the search, and a thread that reads the size of $B_{d,t}$ after $\tau_d$ reads at least $k$.

A call of \textsc{Expand}$(N, u)$ \emph{expands $u$ completely} if it examines every neighbor of $u$ and then sets $\mathit{vis}[u]$ (line~\ref{ln:td-vis}). A call that returns has either expanded $u$ completely or read $\mathit{vis}[u]$ set by a call that did. A block is \emph{scanned in level $d$} once a call of \textsc{ScanBlock} for it by a thread in level $d$ has returned.

\subsection{Buckets, Levels, and Plans}
\label{sec:pf-structures}

\begin{lemma}[Buckets]
\label{lem:buckets}
If a thread reads $k$ as the size of a bucket $B$ and afterwards reads $B.\mathit{array}$, then the array it finds holds the first $k$ vertices appended to $B$ at positions $0$ to $k-1$, and these positions hold the same vertices in every array of $B$ for the rest of the search.
\end{lemma}
\begin{proof}
Only the owner of $B$ writes it. It writes each vertex before increasing $\mathit{size}$ (line~\ref{ln:bk-size}), and publishes a new array only after copying the vertices appended so far into it (line~\ref{ln:bk-array}). When the size became $k$, the array then in $\mathit{array}$ held the first $k$ vertices, and every array published later holds copies of them. The reader reads $\mathit{array}$ after the size, so it finds that array or a later one, and in either case positions $0$ to $k-1$ hold the first $k$ vertices. No array is freed during the search, and no position of an array is written twice.
\end{proof}

\begin{lemma}[The level list]
\label{lem:list}
Every node has at most one successor, which is linked by a CAS from null and never replaced, and a node is fully set up, with its depth set, its buckets empty and its $\mathit{planner}$ empty, before it is linked. The nodes linked from $N_0$ thus form a single list, in which the node $d$ steps from $N_0$ has depth $d$, and every thread walks this list from $N_0$, in order.
\end{lemma}
\begin{proof}
A successor is linked only by the CAS of line~\ref{ln:td-cas}, from null, so at most one CAS on the field $\mathit{next}$ of a node succeeds, and nothing writes the field afterwards. A thread sets up its node before trying the CAS. A node whose CAS fails has been seen by no other thread, since no pointer to it was ever published, and its thread changes only its depth before trying it again at a later level. A node linked to a node of depth $d$ has depth $d+1$ (\textsc{AddLevel}). Every thread starts at $N_0$ and moves on only along $\mathit{next}$.
\end{proof}

\begin{lemma}[Plans]
\label{lem:plans}
For every node $N_d$ that some thread enters, every call of \textsc{ChoosePlan}$(N_d)$ returns the same plan, $P_d$, and every thread calls \textsc{ChoosePlan}$(N_d)$ with $P_{d-1}$, where $P_{-1}$ is the plan before level 0.
\end{lemma}
\begin{proof}
Only CASes from empty write $N_d.\mathit{planner}$ (line~\ref{ln:do-claim}), so the first one succeeds, and $\mathit{planner}$ never changes after it. The thread $w$ whose CAS succeeds wrote its proposal before its CAS (line~\ref{ln:do-propose}), and writes it only once, since it calls \textsc{ChoosePlan}$(N_d)$ once. Every call reads $\mathit{planner}$ set, and then returns the proposal of the thread it names (line~\ref{ln:do-adopt}), which is that of $w$. The second claim follows by induction on $d$, since a thread calls \textsc{ChoosePlan}$(N_{d+1})$ with the plan that its call for $N_d$ returned.
\end{proof}

We write $c_d$ for the code of $P_d$, which is 0 if $P_d$ is top-down. Following Section~\ref{sec:do-codes}, the discoveries of a bottom-up level $d$ write the code \textsc{NextCode}$(c_d)$, which is $c_d + 1$ if $0 < c_d < 255$, and 0 otherwise, and its test is \textsc{InFrontier}$(u, d, c_d)$, which reads $\mathit{code}[u] = c_d$ if $c_d > 0$, and $\mathit{dist}[u] = d$ otherwise. A level $e$ \emph{uses} the code $c_e$ if it is positive, and the code that the discoveries of level $e-1$ write if that is positive. A bottom-up plan $P_e$ takes the code that level $e-1$ writes for its discoveries, if that is positive, and otherwise the smallest code larger than every code used by an earlier level, or 0 if that code would exceed 255.

\begin{lemma}[Codes]
\label{lem:codes}
Every level uses at most one positive code, and no two levels use the same one.
\end{lemma}
\begin{proof}
If level $e-1$ writes a positive code for its discoveries and $P_e$ is bottom-up, $P_e$ takes that code, so level $e$ uses only one. We show by induction on $e$ that the code used by level $e$, if any, is larger than every code used by an earlier level. If the code is the one that level $e-1$ writes, it is $c_{e-1} + 1$, and $c_{e-1}$ is the largest code used before level $e$, by induction. Otherwise, it is a code $c_e$ chosen larger than every code used before.
\end{proof}

\subsection{Safety}
\label{sec:pf-safety}

The proof rests on two properties. The first is that every write is \emph{sound}:
\begin{itemize}
    \item[(S1)] every value $x \geq 0$ written to $\mathit{dist}[v]$ is $\delta(v)$;
    \item[(S2)] every vertex appended to a bucket of $N_d$ is in $L_d$;
    \item[(S3)] every value written to $\mathit{code}[v]$ is the value that the plans give level $\delta(v)$: 0, or the code that level $\delta(v)$ uses.
\end{itemize}
The second is that every level is \emph{ready} by the time its first plan is claimed: for every node $N_d$ that some thread enters, at time $\tau_d$,
\begin{itemize}
    \item[(R1)] $\mathit{dist}[u] = \delta(u)$ for every vertex $u$ of $L_0 \cup \cdots \cup L_d$;
    \item[(R2)] every vertex of $L_d$ is in the reference prefix $S_{d,t}$ of some thread $t$;
    \item[(R3)] if levels $d-1$ and $d$ are both bottom-up and $c_d > 0$, then $\mathit{code}[u] = c_d$ for every vertex $u$ of $L_d$.
\end{itemize}
Sound writes never change a distance or a code once it is set, since all writes to $\mathit{dist}[v]$ write $\delta(v)$, and all writes to $\mathit{code}[v]$ write the same value by (S3). So as long as every write is sound, (R1) to (R3) keep holding after $\tau_d$; (R2) does too, by Lemma~\ref{lem:buckets}. Lemmas~\ref{lem:discover} to~\ref{lem:bu-complete} assume that every write made so far is sound, and that (R1) to (R3) hold for $N_d$; Lemma~\ref{lem:ready} then proves both properties together.

\begin{lemma}[Discoveries]
\label{lem:discover}
If a thread in level $d$ appends a vertex $v$ to a bucket of $N_{d+1}$, then $\delta(v) = d+1$, and the values it writes to $\mathit{dist}[v]$ and $\mathit{code}[v]$ are $d+1$ and the value the plans give level $d+1$. If it writes $\mathit{code}[v] = c$ in \textsc{WriteCodes}$(N_d, c)$, then $\delta(v) = d$, and $c$ is the code that level $d$ uses.
\end{lemma}
\begin{proof}
In a top-down level, the thread appends $v$ while expanding a vertex $u$ taken from a bucket of $N_d$, so $\delta(u) = d$ by (S2), and $\delta(v) \leq d + 1$. It appends $v$ only after reading $\mathit{dist}[v] = -1$ in that expansion, after $\tau_d$, when (R1) holds; hence $\delta(v) > d$. Its CAS writes $d+1$ (line~\ref{ln:do-cas}). In a bottom-up level, the thread appends $v$ after reading $\mathit{dist}[v] = -1$ after $\tau_d$, so again $\delta(v) > d$, and after finding a neighbor $u$ of $v$ with $\mathit{code}[u] = c_d > 0$, or, if $c_d = 0$, with $\mathit{dist}[u] = d$. In the first case, some write gave $u$ the code $c_d$, which level $d$ uses, so $\delta(u) = d$ by (S3) and Lemma~\ref{lem:codes}; in the second, $\delta(u) = d$ by (S1). So $\delta(v) = d+1$, and the thread writes $d+1$ and the code that level $d$ writes for its discoveries, which is the value the plans give level $d+1$. In \textsc{WriteCodes}$(N_d, c)$, the vertices come from buckets of $N_d$, and are in $L_d$ by (S2), and $c = c_d$ is the code that level $d$ uses. None of these arguments depends on how long the thread waits between its reads and its writes, since (R1) and the codes keep holding.
\end{proof}

\begin{lemma}[Certificates]
\label{lem:certs}
The writes that let other threads skip work come after that work:
\begin{itemize}
    \item[(a)] When a thread sets $\mathit{vis}[u]$, every neighbor $v$ of $u$ with $\delta(v) = \delta(u) + 1$ has been appended to a bucket of $N_{\delta(v)}$ and has $\mathit{dist}[v] = \delta(v)$.
    \item[(b)] In a top-down level $d$, when thread $o$ sets $\mathit{last}_{d,o} = i$, the vertices at positions $0$ to $i$ of $B_{d,o}$ have been expanded completely, and when a thread sets $\mathit{done}_{d,o}$, every vertex of $S_{d,o}$ has been.
    \item[(c)] In a bottom-up level $d$, when thread $o$ sets $\mathit{last}_{d,o} = k$, or a thread stamps block $k$ of $o$ with $d$, that block has been scanned in level $d$, and when a thread sets $\mathit{done}_{d,o}$, every block of $o$ has been.
    \item[(d)] When thread $o$ sets $\mathit{codesLast}_{d,o} = i$, the vertices at positions $0$ to $i$ of $B_{d,o}$ have the code $c_d$, and when a thread sets $\mathit{codesDone}_{d,o}$, every vertex of $S_{d,o}$ has it.
\end{itemize}
\end{lemma}
\begin{proof}
Each of these writes follows, in the program order of its thread, either the work it certifies, or the reading of an earlier certificate for that work.

(a) A thread sets $\mathit{vis}[u]$ only after examining every neighbor $v$ of $u$. For each $v$, it either read $\mathit{dist}[v] = -1$ and appended $v$ before its CAS on $\mathit{dist}[v]$, which writes $\delta(v)$ or fails because $\mathit{dist}[v]$ is already set; or it read $\mathit{dist}[v]$ set, by a thread that had appended $v$ before, since both \textsc{Discover} and \textsc{ScanBlock} append a vertex before writing its distance. The appends were to a bucket of $N_{\delta(v)}$, by (S2).

(b) Thread $o$ sets $\mathit{last}_{d,o} = i$ after its call of \textsc{Expand} on position $i$ has returned (line~\ref{ln:td-last}), and a call returns only after $u$ has been expanded completely. A thread sets $\mathit{done}_{d,o}$ (line~\ref{ln:td-done}) after reading the size of $B_{d,o}$ after $\tau_d$, so that the positions it read cover $S_{d,o}$, and after calling \textsc{Expand} on every one of them, or, as a helper, on every one after $\mathit{last}_{d,o}$ as it read it, the positions up to which are covered by the first claim.

(c) Thread $o$ publishes $\mathit{last}_{d,o} = k$ (line~\ref{ln:bu-last}) after scanning block $k$, or after reading its stamp equal to $d$; a helper stamps a block (line~\ref{ln:bu-stamp}) only after scanning it. A stamp equal to $d$ was written after a scan in level $d$: stamps start at $-1$ in every search, and a thread in a level other than $d$ stamps with another depth. The owner sets $\mathit{done}_{d,o}$ after its loop over all of its blocks, which it leaves early only on reading $\mathit{done}_{d,o}$ set already. A helper sets it after taking every block after $\mathit{last}_{d,o}$ as it read it, skipping a block only if the block is now at or before $\mathit{last}_{d,o}$, or stamped with $d$, and stopping only on reading $\mathit{done}_{d,o}$ set already.

(d) As in (b), with \textsc{WriteCodes} in place of \textsc{ExpandLevel}.
\end{proof}

\begin{lemma}[Top-down levels]
\label{lem:td-complete}
If $P_d$ is top-down, then when a thread finishes $N_d$, every vertex $v$ of $L_{d+1}$ has been appended to a bucket of $N_{d+1}$ and has $\mathit{dist}[v] = d+1$, and if $L_d$ is not empty, $N_{d+1}$ has been linked.
\end{lemma}
\begin{proof}
Let $X$ be the thread, let $v \in L_{d+1}$, and let $u$ be a neighbor of $v$ in $L_d$. By (R2), $u$ is at some position $i$ of a reference prefix $S_{d,o}$, and $X$ visited entry $o$ of $N_d$ after $\tau_d$. If it found $\mathit{done}_{d,o}$ set, $u$ had been expanded completely, by Lemma~\ref{lem:certs}(b). Otherwise, the size of $B_{d,o}$ that $X$ read was more than $i$, so the bucket was not empty. If $o = X$, $X$ called \textsc{Expand} on position $i$; if not, it either did so or found $\mathit{last}_{d,o} \geq i$, in which case $u$ had been expanded completely, by Lemma~\ref{lem:certs}(b). A call of \textsc{Expand} returns only after $u$ has been expanded completely, so in every case $u$ has been, and by Lemma~\ref{lem:certs}(a), $v$ has been appended to a bucket of $N_{d+1}$ and has $\mathit{dist}[v] = d+1$. Finally, if $L_d$ is not empty, take any $u \in L_d$ and its entry $o$ as above. Either $X$ found $\mathit{done}_{d,o}$ set, by a thread that had linked $N_{d+1}$ or found it linked before it took the bucket (line~\ref{ln:td-add}), or $X$ took the bucket itself and did the same.
\end{proof}

\begin{lemma}[Bottom-up levels]
\label{lem:bu-complete}
If $P_d$ is bottom-up, then when a thread finishes $N_d$, $N_{d+1}$ has been linked, and every vertex $v$ of $L_{d+1}$ has been appended to a bucket of $N_{d+1}$, has $\mathit{dist}[v] = d+1$, and has the code that level $d$ writes for its discoveries.
\end{lemma}
\begin{proof}
Let $X$ be the thread. It linked $N_{d+1}$, or found it linked, before scanning (\textsc{ScanLevel}). Let $v \in L_{d+1}$, in block $k$ of thread $o$. Thread $X$ took every block of its own and then visited every other thread, so by Lemma~\ref{lem:certs}(c), whichever way it passed block $k$ of $o$, whether by scanning it, finding it at or before $\mathit{last}_{d,o}$ or stamped with $d$, or finding $\mathit{done}_{d,o}$ set, the block has been scanned in level $d$ by some thread $Y$. The call of \textsc{ScanBlock} of $Y$ read $\mathit{dist}[v]$ after $\tau_d$. If it read a value other than $-1$, a thread had discovered $v$, writing its code and appending it before its distance. Otherwise, $v$ was a candidate, and $Y$ tested its neighbors, among them a vertex $u$ of $L_d$. If $c_d = 0$, the test reads $\mathit{dist}[u] = d$, which holds by (R1). If $c_d > 0$, the test reads $\mathit{code}[u] = c_d$: by (R3) if level $d-1$ was bottom-up, and otherwise because $Y$ had finished \textsc{WriteCodes}$(N_d, c_d)$ before scanning (line~\ref{ln:bu-codes}), which gave every vertex of $L_d$ its code, by Lemma~\ref{lem:certs}(d) and (R2), as for expansions in Lemma~\ref{lem:td-complete}. Either way the test succeeds, and $Y$ appended $v$, wrote its code, and set $\mathit{dist}[v] = d+1$ (lines~\ref{ln:bu-append} to~\ref{ln:bu-dist}).
\end{proof}

\begin{lemma}[Soundness and readiness]
\label{lem:ready}
Every write is sound, and (R1) to (R3) hold for every node $N_d$ that some thread enters, from $\tau_d$ to the end of the search.
\end{lemma}
\begin{proof}
Suppose some write is not sound, and let $w$ be the first such write. We first show by induction on $d$ that (R1) to (R3) hold for every node $N_d$ with $\tau_d$ before $w$, from $\tau_d$ up to $w$. For $N_0$ they hold before any call starts (Section~\ref{sec:td-structure}). For $d > 0$, the thread whose CAS on $N_d.\mathit{planner}$ came first finished $N_{d-1}$ before $\tau_d$, after $\tau_{d-1}$. By induction, (R1) to (R3) held for $N_{d-1}$ in that time, and every write before $w$ is sound, so by Lemma~\ref{lem:td-complete} or~\ref{lem:bu-complete}, when it finished $N_{d-1}$, every vertex of $L_d$ had been appended to a bucket of $N_d$ and had its distance, and if level $d-1$ was bottom-up, the code that level $d-1$ writes for its discoveries, which is $c_d$ if level $d$ is bottom-up too and $c_d > 0$. The appends came before $\tau_d$, so they are in the reference prefixes, and (R1) for the earlier levels holds by induction. This gives (R1) to (R3) at $\tau_d$, and they keep holding up to $w$, since the writes before $w$ are sound. Now $w$ is made by a thread in some level $d$, after the reads of that thread that it depends on, which came after $\tau_d$. Lemma~\ref{lem:discover} then shows that $w$ is sound, a contradiction. Every write is therefore sound, and the induction above, which then needs no bound $w$, gives (R1) to (R3) for every node that some thread enters.
\end{proof}

\subsection{The End of the Search}
\label{sec:pf-end}

\begin{lemma}[The last level]
\label{lem:end}
The list consists of the nodes $N_0$ to $N_{D+1}$. A thread that finishes $N_d$ for $d \leq D$ finds $N_{d+1}$ linked, and a thread that finishes $N_{D+1}$ returns.
\end{lemma}
\begin{proof}
For $d \leq D$, $L_d$ is not empty, so by Lemmas~\ref{lem:td-complete} and~\ref{lem:bu-complete}, $N_{d+1}$ has been linked when a thread finishes $N_d$, and the thread finds it (line~\ref{ln:td-next}). By (S2), every vertex appended to a bucket of $N_{D+1}$ is in $L_{D+1}$, which is empty, so its buckets stay empty. Every proposal for $N_{D+1}$ therefore finds $n_f = 0$, and $P_{D+1}$ is top-down (Section~\ref{sec:do-end}). In \textsc{ExpandLevel}$(N_{D+1})$, every entry is empty, so no thread calls \textsc{AddLevel}$(N_{D+1})$, the only call that can link a successor to $N_{D+1}$, and a thread that finishes $N_{D+1}$ finds $\mathit{next}$ null and returns.
\end{proof}

\begin{theorem}[Correctness]
\label{thm:correct}
When any call of \textsc{BFS}$(s)$ returns, $\mathit{dist}[v] = \delta(v)$ for every vertex $v$ that $s$ reaches, and $\mathit{dist}[v] = -1$ for every other vertex, and no later step of any thread changes $\mathit{dist}$.
\end{theorem}
\begin{proof}
By Lemma~\ref{lem:end}, a call returns only after finishing $N_{D+1}$, and hence after $\tau_{D+1}$, when (R1) holds for $N_{D+1}$: every vertex of $L_0 \cup \cdots \cup L_D$, which are all the vertices $s$ reaches, has its distance. By (S1), no thread writes a value $x \geq 0$ to $\mathit{dist}[v]$ for a vertex $v$ that $s$ does not reach, so these distances stay $-1$. Every later write to $\mathit{dist}[v]$ writes $\delta(v)$, the value already there.
\end{proof}

The method of Section~\ref{sec:wf-topdown} is the special case in which every plan is top-down, and the same proof applies to it, with $\tau_d$ the time of the first fence that a thread executes before expanding $N_d$ (line~\ref{ln:td-fence}).

\subsection{Wait-Freedom}
\label{sec:pf-waitfree}

\begin{lemma}[Appends]
\label{lem:once}
In a search, a thread appends each vertex at most once, so all the buckets together hold at most $Tn + 1$ vertices.
\end{lemma}
\begin{proof}
A thread appends a vertex $v$ only after reading $\mathit{dist}[v] = -1$, and right after appending it, the thread writes $\mathit{dist}[v]$, or its CAS on it fails because $\mathit{dist}[v]$ is already set. A distance never returns to $-1$, so the thread never reads $\mathit{dist}[v] = -1$ again. The one vertex not appended this way is the source in $N_0$.
\end{proof}

\begin{theorem}[Wait-freedom]
\label{thm:waitfree}
Every call of \textsc{BFS}$(s)$ returns after $O((D+2)(Tn + m) \log n)$ steps of its own thread, whatever the other threads do.
\end{theorem}
\begin{proof}
A call walks the list once, from $N_0$ to $N_{D+1}$ (Lemmas~\ref{lem:list} and~\ref{lem:end}). No loop of the method waits for another thread: every loop runs over a range fixed when it starts, which is the $T$ entries of a node, the positions of a bucket up to the size read, the neighbors of a vertex, or the blocks of a thread, and the only tests of another thread's progress, of a flag $\mathit{done}$, end a loop early. Each CAS is tried once. We count the steps of one call. Choosing a plan reads the $T$ entries of a node and makes at most one CAS, and linking a successor sets up $T$ entries and makes one CAS: $O(T)$ steps per level. Over the whole search, a thread takes every position of every bucket at most once in each of \textsc{ExpandLevel} and \textsc{WriteCodes}, at most $Tn + 1$ positions in all by Lemma~\ref{lem:once}, and it expands each vertex completely at most once, since it then reads its own $\mathit{vis}[u]$, so it examines at most $m$ neighbors in top-down levels in all. In a bottom-up level, it scans each block at most once, reading each of its vertices and the neighbors of the candidates: $O(n + m)$ steps per level. Appends take constant time amortized, and the copies made as buckets grow add up to fewer positions than the buckets' final arrays, $O(Tn)$ in all. Each random order needs a step coprime to its size $k$, which the search for it, moving from a random start and wrapping to 1, finds within $k$ tries of $O(\log k)$ steps each; the orders cost $O((Tn + m) \log n)$ steps over the positions of buckets and the neighbors of vertices, and $O(n \log n)$ steps per level over blocks. Every allocation takes a bounded number of steps (Section~\ref{sec:bg-model}). The sum is $O((D+2)(Tn + m) \log n)$, which does not depend on the steps of any other thread.
\end{proof}

The bound is loose, and says nothing about speed; what matters for wait-freedom is that it is finite and depends only on $n$, $m$, $T$ and $D$, not on the other threads.

\subsection{Total Store Order}
\label{sec:pf-tso}

In our model, memory follows TSO rather than sequential consistency (Section~\ref{sec:bg-model}). We model a TSO execution in the usual way~\cite{x86tso}: a store goes first to its thread's store buffer, and later, in program order, to memory; a load returns the newest store to its location in its thread's buffer if there is one, and the value in memory otherwise; a CAS waits until its thread's buffer is empty, and then acts on memory atomically; and a fence waits until its thread's buffer is empty. We take the time of a store to be the time it reaches memory. Two properties then hold that sequential consistency gives: the stores of a thread reach memory in program order, and its loads and stores take effect in program order, except that a load can take effect before an earlier store of its thread reaches memory.

Every step of the proof draws its conclusions in one of two ways. In the first, a thread reads its own stores, which it always sees. In the second, a thread $X$ reads a value that another thread $Y$ stored, and concludes that it also sees the stores that $Y$ made, or read from other threads, before that store. Three properties of TSO make this hold. The stores of $Y$ reach memory in program order, so its earlier stores are there by the time the one that $X$ read is. A load of $Y$ takes effect before the stores that follow it, so the stores of other threads that $Y$ read had reached memory before its own store did. And the loads of $X$ take effect in program order, so the loads it makes afterwards find all of these stores, or later ones to the same locations.

Only one step needs more. (R2) must hold in memory at $\tau_d$, but the thread that finished $N_{d-1}$ first may still hold some of its appends to the buckets of $N_d$ in its store buffer. This is why $\tau_d$ is the time of the first CAS on $N_d.\mathit{planner}$, or, in the method of Section~\ref{sec:wf-topdown}, of the first fence before $N_d$. The thread that makes this CAS or fence finished $N_{d-1}$ before it, and either one empties its store buffer. At $\tau_d$, memory therefore holds every append that this thread made or saw, which covers all of $L_d$, together with their distances and codes. Every thread reads the buckets of $N_d$ only after its call of \textsc{ChoosePlan}$(N_d)$ returns, or after its own fence, so after $\tau_d$. Every lemma above therefore holds under TSO, with ``has been appended'' and ``has been set'' read as ``has reached memory, or is in the thread's own store buffer''.

\section{Bounding Memory}
\label{sec:variants}

The method of Section~\ref{sec:dir-opt} allocates a node for every level, and keeps every node and every bucket array until the search ends, since a delayed thread may still read them. Its memory therefore grows with the number of levels, and with all the vertices appended during the search: a search of the road network with 64 threads allocates a node of 64 entries for each of its 6,262 levels, about 400,000 entries in all, although nearly all of them are finished long before the search ends. This section describes a variant that uses only two nodes, one for the even levels and one for the odd, and reuses each of them two levels later. A thread's bucket in a node then keeps its array from one level to the next, so the memory of the variant depends on the largest levels rather than on all of them. The difficulty is that a node is reused while threads that have fallen behind may still be working on the level it held before.

\subsection{Reusing Nodes}
\label{sec:mem-reuse}

The two nodes are $M_0$ and $M_1$, and level $d$ uses $M_{d \bmod 2}$. Every bucket carries a \emph{tag}, the level that its vertices belong to. A thread starts level $d$ by emptying its own bucket in $M_{(d+1) \bmod 2}$, which held level $d-1$, and tagging it with $d+1$ (Algorithm~\ref{alg:mem}, line~\ref{ln:mem-tag}); the thread is past level $d-1$ by then. A thread in level $d$ uses a bucket only if it is tagged $d$. A smaller tag means that the owner has not started level $d-1$, and so has appended nothing to the bucket for level $d$. A larger tag means that the owner has started level $d+1$, which it can do only once level $d+1$ has been chosen, by a thread that had finished level $d$.

A thread that found a bucket tagged $d$ can still be overtaken by its owner, which may empty it and refill it for level $d+2$ while the thread reads it. The thread then meets vertices of level $d+2$, or memory not yet written in this search. It therefore takes a vertex $u$ from another thread's bucket only if $u$ is a vertex number and $\mathit{dist}[u] = d$. Every such vertex is in $L_d$, so a thread expands, and writes codes for, only vertices of its own level, whatever it reads.

The fields that record progress change accordingly. Before tagging its bucket in the other node, a thread resets its fields $\mathit{last}$ and $\mathit{codesLast}$ in $M_{d \bmod 2}$, and another thread counts them for level $d$ only while the owner's bucket in the other node is tagged $d+1$. The flag $\mathit{done}$ becomes the latest level at which the owner's work is known to be finished, and a thread in level $d$ skips an owner whose $\mathit{done}$ is at least $d$; a delayed thread may write a smaller level over a larger one, which only makes other threads check finished work again. Stamps already hold levels (Section~\ref{sec:do-bu}).

Instead of a field $\mathit{planner}$ and proposals, each node has a single word that holds the latest level whose plan was chosen, with the direction of that plan. A thread that finds an older level in the word computes a plan and moves the word to its level with a CAS; whether its CAS succeeds or not, it then follows the direction in the word. Besides top-down and bottom-up, the direction can be \emph{done}: the first thread to find all the buckets of a level empty chooses it, and every thread that reads it returns. A thread that finds a later level in the word knows that its own level is over, since a later level has been chosen, by a thread that had finished every level before it. It then moves on to the older of the two levels in the words of the nodes, which is past its own, and continues from there.

Codes follow the depth instead of the plans: level $d$ has the code $d+1$, as long as it fits in a byte. A thread that falls behind skips the plans of the levels it jumps over, so it could not know codes that depend on them. Whether a bottom-up level writes the codes of its frontier first, which depends on the direction of the level before, is chosen together with its direction and kept in the same word, since only the thread that chooses the plan is sure to know the direction of the level before.

\begin{algorithm}
\caption{The levels of the variant with two nodes, $M_0$ and $M_1$, as thread $t$ runs them. A plan word $W$ holds a level, $W.\mathit{level}$, and the direction chosen for it, $W.\mathit{direction}$.}
\label{alg:mem}
\begin{algorithmic}[1]
\Procedure{BFS}{$s$}
    \State $d \gets 0$; $P \gets$ the plan before level $0$
    \Loop
        \State $N \gets M_{d \bmod 2}$; $N' \gets M_{(d+1) \bmod 2}$
        \State $W \gets N.\mathit{word}$
        \If{$W.\mathit{level} < d$}
            \State $W' \gets$ \Call{Rule}{$P$, $N$, $d$} \Comment{level $d$ and a direction}
            \State \Call{CAS}{$N.\mathit{word}$, $W$, $W'$} \label{ln:mem-cas}
            \State $W \gets N.\mathit{word}$
        \EndIf
        \If{$W.\mathit{level} > d$} \Comment{level $d$ is over}
            \State $d \gets \min(M_0.\mathit{word}.\mathit{level}, M_1.\mathit{word}.\mathit{level})$ \label{ln:mem-jump}
            \State \textbf{continue}
        \EndIf
        \IfThen{$W.\mathit{direction}$ is done}{\Return}
        \State $P \gets$ the plan of level $d$, with direction $W.\mathit{direction}$
        \State $N.\mathit{last}[t] \gets -1$; $N.\mathit{codesLast}[t] \gets -1$
        \State empty $N'.B[t]$ and tag it with $d+1$ \label{ln:mem-tag}
        \State expand level $d$ in direction $W.\mathit{direction}$
        \State $d \gets d+1$
    \EndLoop
\EndProcedure
\end{algorithmic}
\end{algorithm}

\subsection{Correctness and Wait-Freedom}
\label{sec:mem-proof}

\begin{theorem}
\label{thm:mem}
Theorems~\ref{thm:correct} and~\ref{thm:waitfree} hold for the variant.
\end{theorem}
\begin{proof}
We follow Section~\ref{sec:proof}, with $N_d$ standing for $M_{d \bmod 2}$ from the first CAS that moves its word to $d$, at time $\tau_d$, and $B_{d,t}$ for the bucket of thread $t$ in it while the bucket is tagged $d$. Lemma~\ref{lem:plans} holds with the word of the node in place of $\mathit{planner}$ and the proposals, since every thread in level $d$ follows the direction that the first CAS wrote, and Lemma~\ref{lem:codes} holds because level $d$ uses the code $d+1$. The proof then needs four additions.

First, no part of $M_{d \bmod 2}$ is reused for level $d+2$ before level $d+1$ has been chosen: a thread empties and retags its bucket there, and resets its progress there, only once it has started level $d+1$ or $d+2$, and a word moves to a level only after the level before it has been chosen. The thread whose CAS first chooses level $d+1$ finished level $d$ before, so throughout its pass over level $d$, it sees $M_{d \bmod 2}$ just as Section~\ref{sec:proof} sees $N_d$, except that the bucket of an owner that has not started level $d-1$ holds an older tag and is empty for level $d$. Lemmas~\ref{lem:certs} to~\ref{lem:bu-complete} therefore hold for this thread, and with them (R1) to (R3) for level $d+1$ at $\tau_{d+1}$, as in Lemma~\ref{lem:ready}.

Second, a thread that falls behind can meet reused buckets, but it takes from them only vertices $u$ with $\mathit{dist}[u]$ equal to its level, so every write it makes is sound, by Lemma~\ref{lem:discover}. The certificates it writes after level $d+1$ has been chosen can claim work at level $d$ that it never did. By then, however, the thread that chose level $d+1$ has finished level $d$, so these certificates can only make other threads still in level $d$ skip work that the search no longer needs. A later level ignores them: a $\mathit{done}$ of $d$ is below its own level, and an owner's progress counts only next to the owner's current tag.

Third, the word of a node becomes done only when the thread that moves it there has finished the level before and found the buckets of its level empty. That level is $L_{D+1}$: a thread that has finished level $d-1$ finds the vertices of $L_d$, if there are any, in the buckets tagged $d$, as in (R2). Every call therefore returns at level $D+1$, after $\tau_{D+1}$, and the proof of Theorem~\ref{thm:correct} applies.

Fourth, the level of a thread only grows, since a jump (line~\ref{ln:mem-jump}) goes to a level that has been chosen and is past its own. A call therefore takes at most $D+2$ levels, each in a bounded number of its own steps as in Theorem~\ref{thm:waitfree}, and a jump takes a constant number. Under TSO, the CAS on the word of a node (line~\ref{ln:mem-cas}) takes the place of the CAS on $\mathit{planner}$ in Section~\ref{sec:pf-tso}.
\end{proof}

\subsection{Memory}
\label{sec:mem-memory}

The variant keeps two nodes of $T$ entries each, and the bucket of a thread in a node keeps its array for every level that the node holds. The array grows to the largest number of vertices that the thread appends to a level of that parity; its old arrays stay until the search ends, but together hold fewer positions than its current one. Since a thread appends a vertex at most once (Lemma~\ref{lem:once}), the variant needs $O(T)$ entries and $O(\sum_t \max_d |B_{d,t}|) \subseteq O(T \max_d |L_d|)$ positions, against $O(T(D+2))$ entries and $O(Tn)$ positions for the method of Section~\ref{sec:dir-opt}. The two are close on low-diameter graphs, whose largest levels hold a large share of all vertices, and far apart on high-diameter graphs such as road networks, which have thousands of small levels. Section~\ref{sec:experiments} compares the times of the two.

\section{Implementation}
\label{sec:impl}

The methods of Sections~\ref{sec:wf-topdown}, \ref{sec:dir-opt} and~\ref{sec:variants}, and their proofs, are stated in the model of Section~\ref{sec:bg-model}, which speaks of loads, stores, CASes and fences of words rather than of any programming language. An implementation keeps the guarantees of the proofs if it realizes these operations as the model describes them: each atomic, in the order that the model keeps, and each in a bounded number of steps of its thread. How to do so depends on the language, the compiler and the processor. This section describes our implementations of the three methods, which are written in C++ and built with GCC for x86-64 (Section~\ref{sec:exp-setup}).

In C++, two accesses to the same location by different threads, at least one of them a write, form a data race unless both are atomic, and a program with a data race has undefined behavior~\cite{boehm_cpp}. A compiler may therefore assume that no other thread changes a location while a thread accesses it with ordinary loads and stores, and transform these accesses in ways that are harmless for one thread but not for several: it may split an access into smaller ones, or load a location again instead of keeping a value that it loaded, so that two uses of one read in the source see different values~\cite{boehm_benign}. Our methods share locations in exactly this way: helpers read buckets, progress fields and flags while their owners write them, and a thread that has fallen behind writes distances, codes and stamps that other threads are reading.

Our implementations therefore make every access to a location that another thread may access at the same time, with one of the two a write, an atomic one. Most are \emph{relaxed} loads and stores, which are atomic but order nothing: those of distances, codes and flags $\mathit{vis}$; of the sizes, arrays and positions of buckets; of the progress fields $\mathit{last}$, $\mathit{done}$, $\mathit{codesLast}$ and $\mathit{codesDone}$ and the edge counts of entries; of stamps; and, in the variant, of the tags of buckets. We make them with \texttt{std::atomic\_ref}, part of C++ since C++20, which applies atomic operations to an ordinary object. The arrays thus stay ordinary arrays, and the code that sets them up and resets them between searches, when no other thread accesses them, stays ordinary code. A relaxed load or store is indivisible, and the compiler neither loads the location again in place of a value already loaded, nor stores to it where the source does not.

Ordinary accesses remain for memory that only one thread uses during a search, such as the capacity of a bucket; for memory that no thread writes during a search, such as the graph; and for memory that a thread writes before publishing it with a CAS, and that other threads read only after reading what the CAS wrote: a node, which its thread sets up before linking it (Lemma~\ref{lem:list}), and a proposal, which its thread writes before claiming the level (Lemma~\ref{lem:plans}). The fields $\mathit{next}$ and $\mathit{planner}$, and the word of a node in the variant, are \texttt{std::atomic} objects, whose loads and CASes are sequentially consistent. In C++, a load that reads the value written by such a CAS orders the accesses that came before the CAS before those that come after the load, so these ordinary accesses are not data races either.

In the variant, a thread that reads a reused bucket can read positions of its array that no thread has written in the search (Section~\ref{sec:mem-reuse}). The model lets it read them, since every word holds some value, but C++ gives memory that a program has never written an indeterminate value, which a program may not read. The arenas of the variant therefore take their memory from the system zeroed, with \texttt{calloc}, so that every position holds 0 or a value written in an earlier search, which \textsc{Valid} discards unless it is a vertex of the reader's level.

The distances are final once the first call of a search returns, but calls that are still running may write them again, with the values already there (Section~\ref{sec:bg-method}). A consumer that reads them before every call has returned must therefore read them with atomic loads too, or first copy them with atomic loads into memory of its own; it finds the final distances either way.

Relaxed accesses can be reordered: the compiler may move them past each other, and a processor whose memory model is weaker than TSO may make them take effect out of order. On x86-64, GCC compiles every relaxed load and store to a single \texttt{mov} instruction, which TSO keeps in program order with the other accesses of its thread, except a store and a later load, so only the compiler can break the orders that the model needs. We keep that order with \emph{compiler barriers}, empty \texttt{asm} statements that clobber memory, which emit no instruction, and across which GCC moves no memory access. We put one wherever the algorithms require one access of a thread to come before another: between every write that lets other threads skip work and the work it covers (Section~\ref{sec:td-order}); between appending a vertex, writing its code, and writing its distance; in \textsc{Append}, between writing a vertex and the new size, and between copying an array and publishing it (Algorithm~\ref{alg:bucket}); between reading the size of a bucket and its array; and, in the variant, between resetting the progress of a thread, emptying its bucket, and tagging it, and between reading the tag of a bucket and the size or the progress that the tag guards (Algorithms~\ref{alg:mem} to~\ref{alg:mem-codes}). With these barriers, every thread makes its loads and stores in the order that the algorithms list them, and TSO keeps that order as Section~\ref{sec:bg-model} describes.

This way of keeping the order relies on GCC and on TSO. Portable C++ would put the order into the accesses themselves, with a release store for every store that must become visible after the earlier loads and stores of its thread, and an acquire load for every load that must take effect before the later loads and stores of its thread~\cite{boehm_cpp}. On x86-64, GCC compiles release stores and acquire loads to the same \texttt{mov} instructions as relaxed ones, so on our machine the two forms make the same instruction for every access. Processors with weaker memory models, such as ARM and POWER, need instructions of their own for release and acquire, and there the barriers alone would not keep the order.

Every CAS of the methods is a \texttt{compare\_exchange\_strong} with the default, sequentially consistent order, and the fence of Algorithm~\ref{alg:wf-td} (line~\ref{ln:td-fence}) is an \texttt{atomic\_thread\_fence} with the same order. GCC compiles the CAS to a \texttt{lock cmpxchg} instruction, and the fence to a \texttt{lock or} of zero into the top of the stack, which changes no data. A locked instruction takes effect only once the earlier stores of its thread are visible~\cite{intelmanual}, and GCC moves no memory access across a sequentially consistent operation, so both behave as Section~\ref{sec:bg-model} assumes. The fence costs one instruction per level.

The proof of wait-freedom counts every load, store, CAS and fence as one step (Section~\ref{sec:bg-model}), so it carries over to an implementation only if each of them finishes in a bounded number of steps of its own thread. C++ implements an atomic operation that the processor cannot perform directly with a lock, which a delayed thread could hold; every atomic type that we use, of one, four or eight bytes, is always lock-free on x86-64 (\texttt{is\_always\_lock\_free}), and each of our atomic operations is a single instruction. A thread's arena hands out blocks by moving a pointer through chunks of at least 4~MiB, which it keeps from one search to the next, and asks the system allocator for a new chunk only when none of its chunks has room left. These calls, which can take a lock, are the only ones outside the assumption of Section~\ref{sec:bg-model}, and a thread makes them only in a search that needs more memory than its arena already holds.

\section{Experiments}
\label{exp_sec}

Our experiments were conducted on \textbf{AMD EPYC 7452 32-Core Processor} with \textbf{256} GB RAM. The code was compiled using \textbf{g++ 7.5.0} without any compiler optimizations. All the times reported are in microseconds. We only show the plots here, the tables are available in the appendix \ref{sec:tables}.

We evaluate the performance of our wait-free BFS algorithm using random graphs generated with varying sizes and densities as well as real world datasets.

We evaluate two versions of the wait-free BFS algorithm, one that relies on hardware guarantees for atomicity and another that explicitly marks shared variables as atomic.

We also compare our BFS algorithm with Bader and Madduri's \cite{bader_bfs} lock-based BFS algorithm.

% Number experiment references independently of section headings.
\newcounter{bfsexperiment}

\refstepcounter{bfsexperiment}
\paragraph{Experiment \thebfsexperiment: Varying Number of Vertices}
\label{bfs:exp1}
We fix the number of threads at 16 and vary the number of vertices in the graph keeping number of edges proportional to the number of vertices (50 times). Figure~\ref{fig:exp1} present the results.

\pgfplotstableread[col sep=space]{
    Vertices   Serial   WaitFree   WaitFreeAtomic   LockBased
    1e5        139344.6     19755.4    22067.2          61038
    2e5        311584.8     33606.8    37080            114281.8
    3e5        488414.8     47269.8    54004            169902.2
    4e5        663065.4     60732.4    70643.4          230551
    5e5        835574.8     74156.2    86906.2          284763.8
    6e5        1015243      86748.6    103275.8         338586.6
}\bfsA

\begin{figure}
    \centering
    \begin{tikzpicture}
    \begin{axis}[
        width=0.45\textwidth,
        height=0.3\textwidth,
        xlabel={Vertices},
        ylabel={Execution Time ($\mu$s)},
        legend pos=north west,
        xtick=data,
        xticklabels={1,2,3,4,5,6},
        ymajorgrids=true,
        grid style=dashed,
        enlarge x limits=0.05,
        tick label style={font=\small},
        label style={font=\small},
        legend style={font=\small},
        legend style={
            font=\scriptsize,
            cells={anchor=west},
            inner sep=2pt,
            column sep=2pt
        }
    ]
    
    \addplot table[x=Vertices, y=Serial] {\bfsA};
    \addlegendentry{Serial}
    
    \addplot table[x=Vertices, y=WaitFree] {\bfsA};
    \addlegendentry{Wait-Free}
    
    \addplot table[x=Vertices, y=WaitFreeAtomic] {\bfsA};
    \addlegendentry{Wait-Free (Atomic)}
    
    \addplot table[x=Vertices, y=LockBased] {\bfsA};
    \addlegendentry{Lock-Based}
    
    \end{axis}
    \end{tikzpicture}
    \caption{Execution time of BFS for varying number of vertices.}
    \Description{Execution time of BFS for varying number of vertices.}
    \label{fig:exp1}
\end{figure}

From Figure \ref{fig:exp1}, we observe that as the number of vertices increases, the wait-free BFS shows a consistent and substantial speedup compared to the existing parallel BFS with the wait-free version achieving up to $11\times$ speedup over the serial BFS and $4\times$ speedup over the existing parallel BFS by using $16$ threads. The version relying on hardware atomicity performs slightly better than the explicitly atomic version, but both variants achieve a significant speedup.

\refstepcounter{bfsexperiment}
\paragraph{Experiment \thebfsexperiment: Varying Graph Density}
\label{bfs:exp2}
We fix the number of vertices at $5e4$, number of threads at 64 and vary the density of the graph. Figure~\ref{fig:exp2} present the results.

\pgfplotstableread[col sep=space]{
    Density   Serial   WaitFree   WaitFreeAtomic   LockBased
    1         247092.4     20563.6    21823            66997.6
    5         1186723      52643      57167.8          195032
    10        2339696.8    89930.4    100727.2         299954.2
    25        5801718.6    192167     213233.6         506012.4
    50        11578280.2   378123.4   431672.8         814425.8
    100       23380493     737423     818978.8         1431971
}\bfsB

\begin{figure}
    \centering
    \begin{tikzpicture}
    \begin{axis}[
        width=0.45\textwidth,
        height=0.3\textwidth,
        xlabel={Density (Percent)},
        ylabel={Execution Time ($\mu$s)},
        legend pos=north west,
        xtick=data,
        xticklabels={1,5,10,25,50,100},
        ymajorgrids=true,
        grid style=dashed,
        enlarge x limits=0.05,
        tick label style={font=\small},
        label style={font=\small},
        legend style={font=\small},
        legend style={
            font=\scriptsize,
            cells={anchor=west},
            inner sep=2pt,
            column sep=2pt
        }
    ]
    
    \addplot table[x=Density, y=Serial] {\bfsB};
    \addlegendentry{Serial}
    
    \addplot table[x=Density, y=WaitFree] {\bfsB};
    \addlegendentry{Wait-Free}
    
    \addplot table[x=Density, y=WaitFreeAtomic] {\bfsB};
    \addlegendentry{Wait-Free (Atomic)}
    
    \addplot table[x=Density, y=LockBased] {\bfsB};
    \addlegendentry{Lock-Based}
    
    \end{axis}
    \end{tikzpicture}
    \caption{Execution time of BFS algorithms for varying graph densities.}
    \Description{Execution time of BFS algorithms for varying graph densities.}
    \label{fig:exp2}
\end{figure}

From Figure \ref{fig:exp2}, we observe that  the wait-free BFS maintains its efficiency as graph density increases. The speedup remains consistent across different densities, with the wait-free version offering up to $32\times$ faster execution than the sequential BFS and $2\times$ faster than existing parallel BFS using $64$ threads, especially for denser graphs. While both versions exhibit increasing overhead with higher densities, the explicitly atomic variant incurs a modest performance penalty.

\refstepcounter{bfsexperiment}
\paragraph{Experiment \thebfsexperiment: Varying Number of Threads}
\label{bfs:exp3}
We perform the experiment on a random graph with $5e5$ vertices and $2.5e7$ edges. We also perform the experiment two real world datasets, \textit{LiveJournal} \cite{livejournal} and \textit{Sina Weibo} \cite{sinaweibo}. Figure~\ref{fig:exp3_1} present the results for random graph, Figure~\ref{fig:exp3_2} present the results for \textit{LiveJournal} \cite{livejournal} graph and, Figure~\ref{fig:exp3_3} present the results for \textit{Sina Weibo} \cite{sinaweibo} graph.

\pgfplotstableread[col sep=space]{
    Threads   WaitFree   WaitFreeAtomic   LockBased
    2         471627.4     550741      2267025.2
    4         251506.2     296599.6    1036392.2
    8         136896.6     161544.2    553816.8
    16        74470.8      86648.6     285779.8
    32        48459.8      53259.2     217453.2
    64        48378.6      53296       147583.6
}\bfsC

\begin{figure}
    \centering
    \begin{tikzpicture}
    \begin{axis}[
        width=0.45\textwidth,
        height=0.3\textwidth,
        xlabel={Threads},
        ylabel={Execution Time ($\mu$s)},
        legend pos=north east,
        xtick=data,
        xticklabels={2,4,8,16,32,64},
        ymajorgrids=true,
        grid style=dashed,
        enlarge x limits=0.05,
        tick label style={font=\small},
        label style={font=\small},
        legend style={font=\small},
        legend style={
            font=\scriptsize,
            cells={anchor=west},
            inner sep=2pt,
            column sep=2pt
        }
    ]
    
    \addplot table[x=Threads, y=WaitFree] {\bfsC};
    \addlegendentry{Wait-Free}
    
    \addplot table[x=Threads, y=WaitFreeAtomic] {\bfsC};
    \addlegendentry{Wait-Free (Atomic)}
    
    \addplot table[x=Threads, y=LockBased] {\bfsC};
    \addlegendentry{Lock-Based}
    
    \end{axis}
    \end{tikzpicture}
    \caption{Execution time of BFS algorihtms for varying number of threads on random graph.}
    \Description{Execution time of BFS algorihtms for varying number of threads on random graph.}
    \label{fig:exp3_1}
\end{figure}

\pgfplotstableread[col sep=space]{
    Threads   WaitFree   WaitFreeAtomic   LockBased
    2         2126836.8     2650480.2    4365973.8
    4         1130817.4     1420545.4    2232363
    8         601362        755491.6     1192019.2
    16        362708.6      440648.8     623655.4
    32        235820.6      269026.2     399018.6
    64        172725.2      233587.4     294441.6
}\bfsD

\begin{figure}
    \centering
    \begin{tikzpicture}
    \begin{axis}[
        width=0.45\textwidth,
        height=0.3\textwidth,
        xlabel={Threads},
        ylabel={Execution Time ($\mu$s)},
        legend pos=north east,
        xtick=data,
        xticklabels={2,4,8,16,32,64},
        ymajorgrids=true,
        grid style=dashed,
        enlarge x limits=0.05,
        tick label style={font=\small},
        label style={font=\small},
        legend style={font=\small},
        legend style={
            font=\scriptsize,
            cells={anchor=west},
            inner sep=2pt,
            column sep=2pt
        }
    ]
    
    \addplot table[x=Threads, y=WaitFree] {\bfsD};
    \addlegendentry{Wait-Free}
    
    \addplot table[x=Threads, y=WaitFreeAtomic] {\bfsD};
    \addlegendentry{Wait-Free (Atomic)}
    
    \addplot table[x=Threads, y=LockBased] {\bfsD};
    \addlegendentry{Lock-Based}
    
    \end{axis}
    \end{tikzpicture}
    \caption{Execution time of BFS algorithms for varying number of threads on \textit{LiveJournal} \cite{livejournal} graph.}
    \Description{Execution time of BFS algorithms for varying number of threads on \textit{LiveJournal} \cite{livejournal} graph.}
    \label{fig:exp3_2}
\end{figure}

\pgfplotstableread[col sep=space]{
    Threads   WaitFree   WaitFreeAtomic   LockBased
    2         25707716.4     35663801    50833290.2
    4         13594550.8     19191065.4  33334874.2
    8         7132199        10601925    16136456.4
    16        5038638        7583175     10040580.6
    32        2527066.8      4045410.6   6730842.4
    64        2008005        2976349     5643021.2
}\bfsE

\begin{figure}
    \centering
    \begin{tikzpicture}
    \begin{axis}[
        width=0.45\textwidth,
        height=0.3\textwidth,
        xlabel={Threads},
        ylabel={Execution Time ($\mu$s)},
        legend pos=north east,
        xtick=data,
        xticklabels={2,4,8,16,32,64},
        ymajorgrids=true,
        grid style=dashed,
        enlarge x limits=0.05,
        tick label style={font=\small},
        label style={font=\small},
        legend style={font=\small},
        legend style={
            font=\scriptsize,
            cells={anchor=west},
            inner sep=2pt,
            column sep=2pt
        }
    ]
    
    \addplot table[x=Threads, y=WaitFree] {\bfsE};
    \addlegendentry{Wait-Free}
    
    \addplot table[x=Threads, y=WaitFreeAtomic] {\bfsE};
    \addlegendentry{Wait-Free (Atomic)}
    
    \addplot table[x=Threads, y=LockBased] {\bfsE};
    \addlegendentry{Lock-Based}
    
    \end{axis}
    \end{tikzpicture}
    \caption{Execution time of BFS for varying number of threads on \textit{Sina Weibo} \cite{sinaweibo} graph.}
    \Description{Execution time of BFS for varying number of threads on \textit{Sina Weibo} \cite{sinaweibo} graph.}
    \label{fig:exp3_3}
\end{figure}

This experiment explores the performance impact of varying thread counts on both random and real-world graphs. On both the \textit{LiveJournal} and \textit{Sina Weibo} graphs, the wait-free BFS reduces execution time by a factor of $2\times$ to $3\times$ compared to the existing parallel BFS and more than $20\times$ compared to the serial BFS on $64$ threads. The performance improvement is more pronounced in larger real-world graphs, where the algorithm's ability to distribute work efficiently across threads becomes more critical.

\paragraph{Hardware Atomicity vs. Explicit Atomic Operations}
Across all experiments, the variant that relies on hardware atomicity consistently outperforms the explicitly atomic version, albeit by a small margin. The overhead introduced by explicit atomic operations becomes more apparent in high-density graphs and larger thread counts, though the overall performance remains substantially better than the existing BFS.

\paragraph{Summary}
Our experimental evaluation demonstrates that the proposed wait-free BFS algorithm significantly outperforms the existing BFS algorithm across various graph sizes, densities, and thread counts. The wait-free BFS algorithm demonstrates robust scalability and efficiency across diverse experimental settings. It offers substantial performance improvements over existing BFS implementations, making it well-suited for large-scale graph processing in parallel environments.

\section{Related Work}
\label{sec:related}

Parallel breadth-first search (BFS) has been studied extensively on shared-memory systems. A common approach is level-synchronous traversal, which processes the current frontier in parallel before proceeding to the next level. Bader and Madduri~\cite{bader_bfs} developed a fine-grained multithreaded BFS implementation for the Cray MTA-2. Leiserson and Schardl~\cite{charles_bfs} introduced a work-efficient parallel BFS using a bag data structure and Cilk++ reducers to construct and process frontiers through divide-and-conquer parallelism. These approaches established techniques for distributing traversal work and managing frontiers efficiently.

Beamer et al.~\cite{beamer_dobfs} introduced direction-optimizing BFS, which switches between top-down and bottom-up traversal. Top-down traversal examines edges from vertices in the current frontier, whereas bottom-up traversal searches the incoming neighbors of unvisited vertices and stops when a frontier neighbor is found. Switching directions can substantially reduce edge examinations when the frontier is large. The GAP Benchmark Suite (GAPBS)~\cite{gapbs} includes an optimized direction-optimizing BFS implementation, which we use as an experimental baseline. Our direction-optimizing implementation builds on this traversal strategy while changing how threads coordinate and complete traversal work.

Ligra~\cite{ligra} provides a shared-memory graph-processing framework with sparse and dense vertex subsets and parallel operations over vertices and edges. Its adaptive traversal mechanisms support efficient frontier processing across different frontier sizes. Building on Ligra and related infrastructure, the Graph Based Benchmark Suite (GBBS)~\cite{gbbs} combines theoretically efficient parallel graph algorithms with optimized shared-memory implementations. Its frontier-based BFS provides another baseline for our evaluation. Work efficiency and efficient frontier processing are distinct from the progress guarantees considered in this paper.

PASGAL~\cite{pasgal} specifically addresses synchronization and scheduling overhead on large-diameter graphs. It adopts vertical granularity control and hash bags from the reachability-based strong-connectivity algorithm of Wang et al.~\cite{vgc_hashbag}. Vertical granularity control lets a task explore multiple hops locally, reducing the number of global rounds. For BFS, PASGAL combines these local searches with distance updates and multiple frontiers; vertices may be revisited when shorter distances are discovered. Hash bags support frontier construction. PASGAL uses fork-join parallelism, and its reduction of synchronization overhead does not establish a wait-free progress guarantee.

Other parallel BFS methods relax the order between levels. On GPUs, XBFS~\cite{xbfs} makes the bottom-up step asynchronous, so that one iteration can discover vertices of more than one level, but it still proceeds in iterations, each of which starts once the one before has ended. Unordered BFS gives up levels altogether, removing the barriers between them at the cost of potentially more work; Holst and Selin~\cite{holst_unordered} schedule an unordered BFS with a relaxed FIFO queue, and report speedups of 1.6 to 2.4 over a level-synchronous BFS on sparse, high-diameter graphs. Neither makes a search wait-free. XBFS waits for its slowest thread at the end of every iteration, and in an unordered BFS without helping, a thread that is suspended after taking a vertex from the queue holds work that no other thread will do, so that the search cannot complete until it resumes.

Wait-freedom requires every call of a method to complete within a finite number of its own steps, independently of the execution rates of other threads~\cite{herlihy_waitfree}. Concurrent data structures usually achieve it through helping, in which a thread completes the pending operations of others~\cite{aomp, kogan_fastwf}. Wait-freedom has also been studied for concurrent graph data structures. Kallimanis and Kanellou~\cite{kallimanis_graph} provide wait-free graph objects with consistent dynamic traversals, using helping to complete pending operations. Bhardwaj et al.~\cite{graph_snapshots} use multiversioning to take wait-free snapshots of a lock-free dynamic graph, on which queries such as BFS can run during concurrent updates. These works address consistent operations and traversals on a changing graph. Our setting concerns multiple threads cooperating to complete a single BFS on a graph that remains fixed during traversal.

Our contribution is a helping-based execution mechanism instantiated in both top-down and direction-optimizing BFS. Threads can complete outstanding traversal work on behalf of delayed threads, removing blocking waits at level barriers. Under the assumptions of our model (Section~\ref{sec:background}), the resulting BFS methods are wait-free. This contribution complements optimizations to traversal direction, frontier representation, and load balancing. We evaluate the resulting implementations against GAPBS, GBBS, and PASGAL to assess their performance alongside established shared-memory BFS implementations.

\section{Conclusion}
\label{sec:conclusion}

We presented a wait-free method for BFS, in which no thread ever waits for another. The levels of a search form a linked list, with one bucket per thread in every level; a thread expands its own bucket and then helps with what is left of the others, and every write that lets other threads skip work comes after the work it covers, so that a thread that stops at any point leaves nothing that the others cannot redo. Direction optimization keeps this property: the threads agree on the plan of each level with a single compare-and-swap, bottom-up levels divide the vertices into blocks that are owned and helped like buckets, and one-byte codes give the bottom-up step an exact test of the current level, even for a thread that has fallen behind. We proved that every vertex receives its BFS distance, that the distances are final once any call returns, and that every call returns after $O((D+2)(Tn+m)\log n)$ steps of its own thread, under the TSO memory model of x86. A variant that reuses two level nodes keeps these guarantees, and needs memory that depends on the largest levels of a search rather than on all of them.

The method is also fast when no thread is delayed. It is the fastest of the programs we compared on four of the five graphs, 1.33 to 1.89$\times$ faster than the fastest of GAPBS, GBBS and PASGAL, and 1.33 to 2.26$\times$ faster than GAPBS on all five; on the road network, the memory optimized variant is 1.09$\times$ faster than PASGAL. It also keeps its speed when threads stall: on LiveJournal, with one of 16 threads stalled for 10~ms at the start of every level or round, it takes 16\% longer than without stalls, while GAPBS, GBBS and PASGAL take 2.8 to 9.0$\times$ as long, and with 15 of the 16 threads stalled for 1~ms it is still faster than all three (Section~\ref{sec:exp-stalls}).

The road network also shows where the method is weakest. With thousands of thin levels, the cost of sharing a level outweighs its work, and past 24 threads on one socket, more threads make a level slower (Section~\ref{sec:exp-road}). Two ways to reduce that cost are to let threads explore several levels ahead from local queues, as PASGAL does, without giving up wait-freedom, and to make helping prefer work whose memory is close, which matters most across sockets. Our method relies on the order that TSO keeps between the stores of a thread, between its loads, and between a load and a later store; on weaker memory models, such as those of ARM and POWER, the accesses that must stay ordered need release and acquire semantics, which on x86 compile to the same instructions as our accesses but on those processors take instructions of their own (Section~\ref{sec:impl}), and their cost there remains to be measured. The state of a search, its arrays and arenas, can be reset or reclaimed only once every call of the search has returned, so a thread that stops for good keeps that state from ever being reused, as such a thread keeps memory from being reclaimed in other nonblocking methods; neither the result nor the other threads, which can go on to searches with state of their own, wait for it. Finally, nothing in the structure of the method, a list of levels whose per-thread buckets any thread can finish, is specific to BFS, and it could make other frontier-based graph algorithms wait-free, such as connectivity, or single-source shortest paths, whose buckets of tentative distances play the role of levels.


\bibliographystyle{ACM-Reference-Format}
\bibliography{references}

\balance
\appendix
\section*{APPENDIX}

\section{Tables}
\label{sec:tables}

The following tables present the execution times of the BFS algorithms for various configurations, including different numbers of vertices, graph densities, and thread counts. The results are based on the experiments described in the main text.

\begin{table}[H]
    \small
    \setlength{\tabcolsep}{3pt}
    \centering
    \begin{tabular}{|c|c|c|c|c|}
        \hline
        \textbf{Vertices} & \textbf{Serial} & \textbf{Wait-Free} & \textbf{Wait-Free (Atomic)} & \textbf{Lock-Based} \\
        \hline
        $1e5$ & 139344 & 19755 & 22067 & 61038 \\
        $2e5$ & 311584 & 33606 & 37080 & 114281 \\
        $3e5$ & 488414 & 47269 & 54004 & 169902 \\
        $4e5$ & 663065 & 60732 & 70643 & 230551 \\
        $5e5$ & 835574 & 74156 & 86906 & 284763 \\
        $6e5$ & 1015243 & 86748 & 103275 & 338586 \\
        \hline
    \end{tabular}
    \caption{Experiment \ref{bfs:exp1} : Execution time of BFS for varying number of vertices.}
    \label{tab:exp1}
\end{table}

\begin{table}[H]
    \small
    \setlength{\tabcolsep}{4pt}
    \centering
    \begin{tabular}{|c|c|c|c|}
        \hline
        \textbf{Threads} & \textbf{Wait-Free} & \textbf{Wait-Free (Atomic)} & \textbf{Lock-Based} \\
        \hline
        $2$ & 471627 & 550741 & 2267025 \\
        $4$ & 251506 & 296599 & 1036392 \\
        $8$ & 136896 & 161544 & 553816 \\
        $16$ & 74470 & 86648 & 285779 \\
        $32$ & 48459 & 53259 & 217453 \\
        $64$ & 48378 & 53296 & 147583 \\
        \hline
    \end{tabular}
    \caption{Experiment \ref{bfs:exp3} : Execution time of BFS algorihtms for varying number of threads on random graph.}
    \label{tab:exp3_1}
\end{table}

\begin{table}[H]
    \small
    \setlength{\tabcolsep}{4pt}
    \centering
    \begin{tabular}{|c|c|c|c|}
        \hline
        \textbf{Threads} & \textbf{Wait-Free} & \textbf{Wait-Free (Atomic)} & \textbf{Lock-Based} \\
        \hline
        $2$ & 2126836 & 2650480 & 4365973 \\
        $4$ & 1130817 & 1420545 & 2232363 \\
        $8$ & 601362 & 755491 & 1192019 \\
        $16$ & 362708 & 440648 & 623655 \\
        $32$ & 235820 & 269026 & 399018 \\
        $64$ & 172725 & 233587 & 294441 \\
        \hline
    \end{tabular}
    \caption{Experiment \ref{bfs:exp3} : Execution time of BFS algorithms for varying number of threads on \textit{LiveJournal} \cite{livejournal} graph.}
    \label{tab:exp3_2}
\end{table}

\begin{table}[H]
    \small
    \setlength{\tabcolsep}{4pt}
    \centering
    \begin{tabular}{|c|c|c|c|}
        \hline
        \textbf{Threads} & \textbf{Wait-Free} & \textbf{Wait-Free (Atomic)} & \textbf{Lock-Based} \\
        \hline
        $2$ & 25707716 & 35663801 & 50833290 \\
        $4$ & 13594550 & 19191065 & 33334874 \\
        $8$ & 7132199 & 10601925 & 16136456 \\
        $16$ & 5038638 & 7583175 & 10040580 \\
        $32$ & 2527066 & 4045410 & 6730842 \\
        $64$ & 2008005 & 2976349 & 5643021 \\
        \hline
    \end{tabular}
    \caption{Experiment \ref{bfs:exp3} : Execution time of BFS for varying number of threads on \textit{Sina Weibo} \cite{sinaweibo} graph.}
    \label{tab:exp3_3}
\end{table}

\begin{table}[H]
    \small
    \setlength{\tabcolsep}{2.5pt}
    \centering
    \begin{tabular}{|c|c|c|c|c|}
        \hline
        \shortstack{\textbf{Density}\\\textbf{(Percent)}} & \textbf{Serial} & \textbf{Wait-Free} & \textbf{Wait-Free (Atomic)} & \textbf{Lock-Based} \\
        \hline
        $1$ & 247092 & 20563 & 21823 & 66997 \\
        $5$ & 1186723 & 52643 & 57167 & 195032 \\
        $10$ & 2339696 & 89930 & 100727 & 299954 \\
        $25$ & 5801718 & 192167 & 213233 & 506012 \\
        $50$ & 11578280 & 378123 & 431672 & 814425 \\
        $100$ & 23380493 & 737423 & 818978 & 1431971 \\
        \hline
    \end{tabular}
    \caption{Experiment \ref{bfs:exp2} : Execution time of BFS algorithms for varying graph densities.}
    \label{tab:exp2}
\end{table}



\end{document}
